\documentclass[12pt,letterpaper]{article}
\usepackage[T1]{fontenc}
% amsmath and amsthm must be loaded before newtxmath (else \openbox clashes).
\usepackage{amsmath,amsthm,mathtools}
% Times for text and math: newtx when installed, mathptmx otherwise.
\IfFileExists{newtxmath.sty}%
  {\usepackage{newtxtext,newtxmath}}%
  {\usepackage{mathptmx}\usepackage{amssymb}}
\usepackage[letterpaper,margin=1in]{geometry}
\usepackage{booktabs,tabularx,array,enumitem}
\usepackage{microtype}
\usepackage[hidelinks]{hyperref}

\setcounter{tocdepth}{2}
\numberwithin{equation}{section}

\theoremstyle{plain}
\newtheorem{theorem}{Theorem}[section]
\newtheorem{proposition}[theorem]{Proposition}
\newtheorem{lemma}[theorem]{Lemma}
\newtheorem{corollary}[theorem]{Corollary}
\theoremstyle{definition}
\newtheorem{definition}[theorem]{Definition}
\newtheorem{example}[theorem]{Example}
\theoremstyle{remark}
\newtheorem{remark}[theorem]{Remark}
\newtheorem{warning}[theorem]{Caution}
% Statements that were not fully verified in this revision.
\newcounter{uvc}
\renewcommand{\theuvc}{2026.\arabic{uvc}}
\newtheorem{uvremark}[uvc]{Remark}

\newcommand{\Gt}{\widetilde G}
\newcommand{\At}{\widetilde A}
\newcommand{\Bt}{\widetilde B}
\newcommand{\Ct}{\widetilde C}
\newcommand{\hl}{\tilde h}
\newcommand{\wt}{\tilde w}
\newcommand{\Cc}{\mathbb C}
\newcommand{\Zz}{\mathbb Z}
\newcommand{\Qq}{\mathbb Q}
\newcommand{\Of}{\mathcal O}
\newcommand{\ZA}{Z_A}
\newcommand{\gen}{\mathrm{gen}}
\newcommand{\reg}{\mathrm{reg}}
\DeclareMathOperator{\PI}{I}
\DeclareMathOperator{\Ind}{Ind}
\DeclareMathOperator{\tr}{tr}
\DeclareMathOperator{\Hom}{Hom}
\DeclareMathOperator{\vol}{vol}
\DeclareMathOperator{\Ad}{Ad}
\newcommand{\SL}{\mathrm{SL}}
\newcommand{\GL}{\mathrm{GL}}

\title{On the unitary dual of $\overline{SL(2,F)}$, $F$ $p$-adic\\[6pt]
\large Draft 3. Corrected and expanded version}
\author{David Joyner\thanks{This is a corrected and expanded version of
the author's unpublished manuscript of the same title, dated July~21,
1999, which is cited as a preprint in \cite{Jo01}. Draft~1 of the
revision is dated November~9, 2026, and Draft~2 November~13, 2026.
This Draft~3, which answers a referee report on Draft~2, is dated
November~13, 2026.}\\[2pt]
\small with Argon-Gemini, Sol-ChatGPT and Claude\\[2pt]
\small\texttt{wdjoyner@gmail.com}}
\date{Original draft: 1999. Corrected and expanded draft: 2026.}

\begin{document}
\maketitle

\begin{abstract}
This paper was revised with the help of the AI systems Argon-Gemini,
Sol-ChatGPT and Claude. Their roles are described in the
Acknowledgments.

Let $\Gt$ be Kubota's $n$-fold cover of $\SL(2,F)$, where $F$ is a
$p$-adic field containing the $n$th roots of unity. The preimage $\At$
of the diagonal torus is not abelian when $n>2$. We describe its
genuine irreducible representations by a Stone--von Neumann argument.
We prove that normalized split orbital integrals of genuine test
functions vanish off the center of $\At$, and we determine their image
when $\gcd(p,2n)=1$. We compute the trace and the character of a
genuine principal series representation. The character involves Weyl
conjugation in the cover, not inversion. When $\gcd(p,2n)=1$ and $n$ is
odd or divisible by~$4$, we deduce a character identity between genuine
principal series of $\Gt$ and principal series of $\SL(2,F)$ for test
functions whose split orbital integrals match. The case
$n\equiv2\pmod 4$, which includes the double cover, is excluded from
this identity. These are the paper's contributions. We do not
classify the unitary dual. Instead we compare each unitary-dual claim
of the 1999 manuscript and of the author's 2001 paper with what is
proved here, and we survey the literature on the rest.
\end{abstract}

\tableofcontents

%======================================================================
\section{Introduction}\label{sec:intro}

Kubota's cover of $\SL(2,F)$ has a nonabelian torus, and that one fact
shapes everything below. We state the problem and the contributions,
compare them with the 1999 and 2001 accounts of the unitary dual, and
fix the hypotheses.

Let $F$ be a nonarchimedean local field of characteristic zero, with
residue characteristic~$p$. Fix $n\ge1$ and assume that $F$ contains
the group $\mu_n$ of $n$th roots of unity, of order $n$. Kubota
\cite{Ku67} constructed a central extension
\[
1\longrightarrow\mu_n\longrightarrow\Gt
\stackrel{\rho}{\longrightarrow}G=\SL(2,F)\longrightarrow1 .
\]
The 1999 title writes $\overline{SL(2,F)}$ for $\Gt$. A representation
of $\Gt$ is \emph{genuine} if $\mu_n$ acts on it through a fixed
faithful character. Genuine representations are the local objects of
the theory of metaplectic forms of Kubota, Gelbart, Flicker and
Kazhdan--Patterson (Section~\ref{sec:literature}).

The 1999 manuscript had two aims. The first was to describe the
genuine unitary dual of $\Gt$. The second was to relate genuine
principal series of $\Gt$ to principal series of $G$ by a character
identity. The motivation was a planned proof of Labesse's
multiplicity one conjecture for $\SL(2)$ over a number field, which
needed the case $n=2$.

A basic difficulty is that the preimage $\At$ of the diagonal torus is
not abelian when $n>2$. Its genuine irreducible representations have
dimension $d=n/\gcd(n,2)$, and their characters vanish off the center
of $\At$. For example, when $n=3$ every genuine irreducible
representation of $\At$ has dimension three. This affects both the
orbital integrals and the characters of induced representations.

\paragraph{Contributions.} The paper makes three contributions.
\begin{enumerate}[label=(\Alph*),leftmargin=*]
\item An explicit treatment of the covering torus: its center, its
  genuine irreducible representations and their traces, and the action
  of the Weyl element (Section~\ref{sec:torus}).
\item A normalized image theorem for split orbital integrals of genuine
  test functions (Section~\ref{sec:orbital}).
\item A restricted split-torus trace identity between genuine principal
  series of $\Gt$ and principal series of $\SL(2,F)$
  (Section~\ref{sec:transfer}), built on a direct computation of the
  trace and character of genuine principal series
  (Section~\ref{sec:characters}).
\end{enumerate}

\paragraph{Main results.} In detail, the paper proves the following.
\begin{enumerate}[label=(\roman*),leftmargin=*]
\item Theorem~\ref{thm:torus} describes the genuine irreducible
  representations of $\At$ by their central characters and computes
  their traces.
\item Lemma~\ref{lem:vanishing} and Theorem~\ref{thm:orbital} show that
  normalized split orbital integrals vanish off the center $\ZA$ of
  $\At$, and identify their image.
\item Lemma~\ref{lem:induced-trace}, Corollary~\ref{cor:trace-integral}
  and Theorem~\ref{thm:character} compute the trace and the character
  of a genuine principal series representation.
\item Lemma~\ref{lem:Kfixed} shows that an unramified genuine principal
  series has a one-dimensional space of spherical vectors.
\item Theorem~\ref{thm:transfer} gives the character identity
  \[
  \tr\PI(\chi_\lambda)(\bar f)=
  \tr\Ind_B^G(\eta_\lambda\delta_B^{1/2})(f)
  \]
  for test functions $\bar f$ on $\Gt$ and $f$ on $G$ whose split
  orbital integrals match.
\end{enumerate}
Along the way it proves Kubota's explicit splitting of $\SL(2,\Of)$
(Lemma~\ref{lem:splitting}) and the rank-one Jacquet-module computation
(Lemma~\ref{lem:jacquet}), so that these results do not rest on
references with other hypotheses.

\paragraph{How the results depend on each other.}
\begin{center}
\small
\begin{tabular}{@{}l@{\ $\longrightarrow$\ }l@{\ $\longrightarrow$\ }l@{}}
Lemma~\ref{lem:cocycle} (cocycle) & Lemma~\ref{lem:radical} (center)
 & Theorem~\ref{thm:torus} (torus traces)\\
Lemma~\ref{lem:splitting} and \eqref{eq:iwasawa} & Lemma~\ref{lem:abel}
 (constant term) & Theorem~\ref{thm:orbital} (image)\\
Theorem~\ref{thm:torus} and Lemma~\ref{lem:abel} &
 Lemma~\ref{lem:induced-trace} (trace) &
 Theorem~\ref{thm:character} (character)\\
Corollary~\ref{cor:trace-integral} and Lemma~\ref{lem:norm} &
 Theorem~\ref{thm:orbital} (existence) &
 Theorem~\ref{thm:transfer} (matching)
\end{tabular}
\end{center}

\paragraph{Relation to the 1999 and 2001 unitary-dual accounts.} The
1999 manuscript described the genuine unitary dual by reducing it,
through Jacquet's theorem, to principal series, complementary series,
the representations at the end of the complementary interval, and
supercuspidal representations. The author's 2001 paper \cite{Jo01}
restates these claims in its Section~3 and adds a list of the types of
unitary and tempered representations, for $\gcd(p,2n)=1$. That list is
only as reliable as the claims it rests on. Table~\ref{tab:status}
compares each claim with what is proved here. Here $s_{\mathrm{old}}$ is
the 1999 parameter, and Section~\ref{subsec:complementary} converts it
to the parameter $s$ of this paper.

\begin{table}[ht]
\centering
\small
\renewcommand{\arraystretch}{1.15}
\begin{tabularx}{\linewidth}{@{}>{\raggedright\arraybackslash}p{0.33\linewidth}>{\raggedright\arraybackslash}X@{}}
\toprule
Claim of the 1999 manuscript (repeated in \cite{Jo01}) & Status in this
draft\\
\midrule
(a) Genuine representations of $\At$ come 2-to-1 from characters of
$h(\varpi^{n\Zz}\Of^\times)$. & False as stated. Corrected and proved:
Theorem~\ref{thm:torus}.\\
(b) Split orbital integrals of genuine functions fill out the
Weyl-invariant genuine functions on the center. & Proved:
Theorem~\ref{thm:orbital}, with the normalizations compared in
Remark~\ref{rem:image-jo01}.\\
(c) Unitary principal series: irreducible for $n$ even, and for $n$ odd
reducible exactly when $\chi^{\wt}=\chi\ne1$ (attributed to Moen). &
Irreducibility for Weyl-regular data is proved
(Proposition~\ref{prop:tempered}). The reducible cases are
source checked against \cite{GS16} (Remark~\ref{uv:reducible}).\\
(d) $\PI(\chi)$ is a complementary series for
$0<s_{\mathrm{old}}<1/n$, by extending it to a cover of $\GL(2,F)$ and
applying \cite[\S\S26--27]{FK86}. & Not proved. The extension step was
never proved (Remark~\ref{uv:complementary}).\\
(e) At $s_{\mathrm{old}}=1/n$: a square-integrable subrepresentation and
a non-tempered spherical quotient, unitary by \cite[Thm.~II.2.1]{KP84}.
& Cited only (Remark~\ref{uv:complementary}).\\
(f) The list of all unitary and tempered types (in \cite{Jo01}). & Cited
only. It follows from (c)--(e) and Jacquet's theorem
(Remark~\ref{uv:subquotient}).\\
(g) The principal-series character uses $\tr\chi(a)+\tr(\chi(a)^{-1})$.
& False for $n>2$. Corrected and proved with Weyl conjugation:
Theorem~\ref{thm:character}.\\
(h) The character identity holds for every $n$ with $\gcd(p,2n)=1$. &
Proved only for $n$ odd or $4\mid n$: Theorem~\ref{thm:transfer}.\\
\bottomrule
\end{tabularx}
\caption{The 1999 and 2001 claims and their status in this draft.}
\label{tab:status}
\end{table}

So this draft proves (a), (b), (g) and (h) in corrected form, and the
Weyl-regular part of (c). It neither proves nor refutes (d), (e) and
(f). The title is the 1999 title, kept by the conventions of this
series of revisions. A more accurate description of the present
content would be ``the covering torus and split principal-series
characters for Kubota covers of $\SL(2)$''.

\paragraph{Restriction on $n$.} Theorem~\ref{thm:transfer} assumes
$\gcd(p,2n)=1$ and that $n$ is odd or divisible by $4$. It does not
cover $n\equiv2\pmod4$, and in particular it does not cover the double
cover $n=2$ that motivated the 1999 manuscript.
Remark~\ref{rem:weil} explains why.

\paragraph{Status labels.} Every result is of one of four kinds.
\emph{Proved here} results are theorems, propositions and lemmas with
proofs. The others are stated in remarks numbered ``Remark 2026.$n$'',
each of which opens with its status. \emph{Source checked} means that
the statement was compared with the cited source, at least with its
abstract and introduction. \emph{Cited only} means that the source was
not consulted for this revision. \emph{Conjectural} means that no proof
is known to us. Open problems are collected in
Section~\ref{sec:limitations}.

\paragraph{Standing hypotheses.} Throughout, $|\mu_n(F)|=n$. The
description of the center in Lemma~\ref{lem:radical}, and
Lemma~\ref{lem:vanishing}, need nothing more. From
Section~\ref{subsec:svn} on we assume $\gcd(p,n)=1$. When the splitting
of $\SL(2,\Of)$ is used, that is, in the image theorem and in all trace
computations, we assume $\gcd(p,2n)=1$, as in the 1999 manuscript. The
splitting itself needs only $\gcd(p,n)=1$
(Lemma~\ref{lem:splitting}). We do not assume $\mu_{2n}\subset F$. Each
theorem restates its hypotheses.

\paragraph{Outline.} Section~\ref{sec:literature} reviews the
literature. Section~\ref{sec:background} explains the background, and
Section~\ref{sec:notation} collects the notation.
Sections~\ref{sec:torus}--\ref{sec:transfer} contain the results.
Section~\ref{sec:limitations} summarizes them and lists open
questions. A glossary follows.

%======================================================================
\section{Literature review}\label{sec:literature}

The relevant literature divides at 1999, the date of the original
manuscript. A result
for a cover of $\GL(2,F)$ does not automatically give the analogous
result for the cover of $\SL(2,F)$, so we say which group each work
treats.

\subsection{Work before 1999}

\paragraph{The cover and the double cover.} Kubota \cite{Ku67}
constructed the $n$-fold covering group of $\SL(2,F)$ by an explicit
two-cocycle built from the $n$th Hilbert symbol. This cocycle is the
starting point of Sections~\ref{sec:background} and~\ref{sec:torus}.
Gelbart \cite{Ge76} developed the representation theory and the
spectral theory of the double cover, with the Weil representation as
the central example. Gelbart and Sally \cite{GS75} studied the
intertwining operators of the double cover and the decomposition of
the Weil representation. Their work explains why the exceptional
constituents of induced representations must be treated separately
from the generic principal series.

\paragraph{Principal series of covers of $\SL(2)$.} Moen \cite{Mo}
studied intertwining operators for covering groups of $\SL(2)$ over a
local field. Moen \cite{Mo88} proved the irreducibility of certain
unitary principal series of these covers when $p$ is odd. The 1999
manuscript attributes to Moen the statement that for $\gcd(p,n)=1$ all
unitary genuine principal series are irreducible when $n$ is even, and
that for $n$ odd they are reducible exactly when $\chi^{\wt}=\chi\ne1$.
Goldberg and Szpruch later proved a statement of this kind (see
below). Arit\"urk
\cite{Ar80} determined composition series of principal series of the
threefold cover. For example, he showed that off the line
$\operatorname{Re}s=0$ the genuine principal series of the threefold
cover are irreducible except at $s=\pm1/3$. The present paper does not
reprove these results. Section~\ref{sec:unitary} cites them.

\paragraph{Parabolic induction.} Bernstein and Zelevinsky
\cite{BZ76,BZ77} developed parabolic induction and Jacquet modules for
reductive $p$-adic groups. Bernstein's notes \cite{BD84} remark that
these arguments carry over to finite central extensions. This is the
basis of the Jacquet-module statements in Section~\ref{sec:unitary}.

\paragraph{Covers of $\GL$ and the metaplectic correspondence.}
Flicker \cite{Fl80} studied automorphic forms on covers of $\GL(2)$
and compared them with forms on $\GL(2)$ by matching orbital
integrals. Kazhdan and Patterson \cite{KP84} developed metaplectic
forms on covers of $\GL(r)$, including the principal series and the
exceptional representations. Flicker and Kazhdan \cite{FK86} proved
the metaplectic correspondence between genuine representations of
covers of $\GL(r)$ and representations of $\GL(r)$. Their comparison
of orbital integrals is the model for Section~\ref{sec:transfer}. They
do not, by themselves, give a transfer for the cover of $\SL(2)$.

\paragraph{Supercuspidal representations and the author's related
work.} Manderscheid \cite{Ma84a,Ma84b} described the supercuspidal
representations of $\SL(2)$ and its double cover. Joyner \cite{Jo88}
studied a metaplectic analogue of Kazhdan's endoscopic lifting.
Joyner \cite{Jo98} proved a correspondence between generalized Hecke
algebras of $\SL(2,F)$ and of its $n$-fold cover. It gives a
correspondence of representations of sufficiently large level, not a
classification of all of them. Schultz \cite{Sc98} studied the lifting
of characters between the double cover of $\SL(2,F)$ and
$SO_{1,2}(F)$. That is a $p$-adic analogue of Adams' lifting for real
groups, and it concerns the case $n=2$ excluded from
Theorem~\ref{thm:transfer}.

\subsection{Work from 1999 on}

\paragraph{The author's papers.} Joyner \cite{Jo99} described the
supercuspidal representations of the $n$-fold cover when
$\gcd(p,2n)=1$, following Manderscheid. That paper proves a Frobenius
formula for their characters and a character relation with $\SL(2,F)$
when $n=2$. Joyner \cite{Jo01} studied invariant distributions on
covers of $\GL(r,F)$ and the unitary and tempered dual of the cover of
$\SL(2,F)$ for $\gcd(p,2n)=1$. It cites the 1999 manuscript as a
preprint. Its Section~3 restates the unitary-dual claims of the 1999
manuscript and lists the resulting types of representations.
Table~\ref{tab:status} gives the status of each claim. It also states
the image theorem for split orbital integrals in a different
normalization (Remark~\ref{rem:image-jo01}).

\paragraph{General covering groups.} McNamara \cite{Mc12} developed the
principal series of central extensions of split reductive groups. He
uses a Stone--von Neumann theorem for the covering torus, computes the
Jacquet modules of principal series, and proves a Satake isomorphism.
He assumes that $F$ contains the $2n$th roots of unity, that is,
$2n\mid q-1$, and that $p\nmid 2n$. That is stronger than the
hypothesis $\mu_n\subset F$ used here. For example, $F=\Qq_5$ with
$n=4$ is allowed here but not in \cite{Mc12}. So we use \cite{Mc12}
only for comparison. Section~\ref{sec:torus} is the rank-one analogue
of his torus theory, and Lemma~\ref{lem:jacquet} is proved directly. Li \cite{Li12} developed local harmonic analysis for
finite central covers of connected reductive groups. That work
includes a Plancherel theorem \cite[Thm.~2.6.4]{Li12}, a Paley--Wiener
theorem \cite[Thm.~5.6.3]{Li12} and Kazhdan's density theorem
\cite[Thm.~5.8.10]{Li12}. Tang \cite{Ta17} computed Plancherel measures
of principal series of central extensions of split groups by a
comparison method. Tang also constructed genuine central characters of
the metaplectic torus in the simply laced case. Luo \cite{Lu20} proved
the Knapp--Stein dimension theorem for finite central covers.

\paragraph{The cover of $\SL(2)$.} Goldberg and Szpruch \cite{GS16}
computed the Plancherel measures of genuine principal series of the
$n$-fold cover of $\SL(2,F)$ when $\mu_n\subset F$ and $p\nmid n$. As
an application they found all reducible genuine principal series
induced from unitary data \cite[Thm.~5.2]{GS16}. Such a representation
is reducible if and only if $n$ is odd and the projection of its
central character to the $d$th powers is a nontrivial quadratic
character. Szpruch \cite{Sz19} defined a matrix of local coefficients
for these representations. He does not assume that $p\nmid n$, but he
assumes $4\nmid n$. Szpruch \cite{Sz25} computed the Whittaker spaces of
the two summands of a reducible unitary genuine principal series,
without assuming $p\nmid n$.

\paragraph{How this paper fits.} Sections~\ref{sec:torus},
\ref{sec:orbital}, \ref{sec:characters} and~\ref{sec:transfer} are
elementary and self-contained, apart from the Weyl integration formula
and standard facts about the cover. Section~\ref{sec:unitary} relies
on the works cited above, and its statements should be read with the
hypotheses of those works.

%======================================================================
\section{Background}\label{sec:background}

Everything in the paper rests on an explicit cocycle, explicit Haar
measures and the classical theory for $\SL(2,F)$. We cover the Hilbert
symbol, Kubota's cover and its splitting over $\SL(2,\Of)$, genuine
representations and functions, Haar measures, and the orbital integrals
and principal series of the linear group. Each subsection ends with an
example.

\subsection{The Hilbert symbol}\label{subsec:hilbert}

The cover is built from the $n$th Hilbert symbol. We recall the
properties that we use.

Let $\Of$ be the ring of integers of $F$, $\varpi$ a uniformizer and
$q$ the order of the residue field. Write $|\cdot|$ for the absolute
value with $|\varpi|=q^{-1}$, and $v$ for the valuation. The $n$th
Hilbert symbol is a map $(\ ,\ )_n:F^\times\times F^\times\to\mu_n$. We
identify $\mu_n$ with a subgroup of $\Cc^\times$ by a fixed
isomorphism. The symbol has the following properties
\cite{Ne86}.
\begin{enumerate}[label=(H\arabic*),leftmargin=*]
\item It is multiplicative in each variable.
\item It is antisymmetric: $(a,b)_n(b,a)_n=1$.
\item $(a,-a)_n=1$. Hence $(a,a)_n=(a,-1)_n$.
\item It is nondegenerate: $(a,b)_n=1$ for all $b$ if and only if
  $a\in(F^\times)^n$.
\item (Tame formula.) If $p\nmid n$, write $\alpha=v(a)$ and
  $\beta=v(b)$. Then $(a,b)_n$ is the image in $\mu_n$ of
  \[
  \bigl((-1)^{\alpha\beta}a^\beta b^{-\alpha}\bmod\varpi\bigr)^{(q-1)/n}.
  \]
  In particular $(u,u')_n=1$ for units $u,u'$.
\end{enumerate}
Some authors use the inverse of this symbol. Every statement below is
unchanged by that choice.

By (H3), $(t,-1)_n^2=(t,1)_n=1$, so $(t,-1)_n=\pm1$. If $n$ is odd,
then $(t,-1)_n=(t,(-1)^n)_n=(t,-1)_n^n=1$.

\begin{example}\label{ex:hilbert}
Take $F=\Qq_7$ and $n=3$. Then $q=7$ and $\mu_3\subset\Qq_7$. By (H5),
$(7,3)_3$ is the image of $(3^{-1})^2\equiv5^2\equiv4\pmod7$. Since
$4^3\equiv1$ and $4\ne1$, $(7,3)_3$ is a nontrivial cube root of
unity. On the other hand $(2,3)_3=1$, since $2$ and $3$ are units.

Take $F=\Qq_5$ and $n=4$. By (H5), $(5,-1)_4$ is the image of
$(-1)^{-1}\equiv-1\pmod5$, raised to the power $(5-1)/4=1$. So $(5,-1)_4=-1$. This value is used
in Example~\ref{ex:inverse}.
\end{example}

\subsection{Kubota's cover}\label{subsec:kubota}

Kubota's cocycle is explicit enough to compute products of lifts of
diagonal, unipotent and Weyl elements, and to write down a splitting
over $\SL(2,\Of)$.

Let $G=\SL(2,F)$ and put
\[
h(t)=\begin{pmatrix}t&0\\0&t^{-1}\end{pmatrix},\qquad
u(x)=\begin{pmatrix}1&x\\0&1\end{pmatrix},\qquad
w=\begin{pmatrix}0&1\\-1&0\end{pmatrix}.
\]
Let $A=h(F^\times)$, $N=u(F)$, $B=AN$ and $W=N_G(A)/A=\{1,w\}$. For
$g=\left(\begin{smallmatrix}a&b\\c&d\end{smallmatrix}\right)$ put
$x(g)=c$ if $c\ne0$ and $x(g)=d$ if $c=0$. Kubota's cocycle is
\begin{equation}\label{eq:kubota}
\beta(g_1,g_2)=
\left(\frac{x(g_1g_2)}{x(g_1)},\frac{x(g_1g_2)}{x(g_2)}\right)_n .
\end{equation}
It is a two-cocycle \cite[\S1]{Ku67}. The group $\Gt$ is the set
$G\times\mu_n$ with the product
$(g,\zeta)(g',\zeta')=(gg',\zeta\zeta'\beta(g,g'))$, and $\rho$ is the
projection to $G$. For $H\subset G$ write $\widetilde H=\rho^{-1}(H)$.
Put $\At=\widetilde A$, $\Bt=\widetilde B$,
\[
\hl(t)=(h(t),1),\qquad \wt=(w,1).
\]
We identify $\zeta\in\mu_n$ with $(1,\zeta)\in\Gt$, which is central.

\begin{lemma}\label{lem:cocycle}
\begin{enumerate}[label=(\alph*)]
\item $\beta(h(a),h(b))=(b,a)_n$.
\item $\beta(u(x),u(y))=1$ and $\beta(h(t),u(x))=1$. Hence
  $u(x)\mapsto(u(x),1)$ is a splitting of $N$, and
  $\hl(t)(u(x),1)=(h(t)u(x),1)$.
\item The splitting of $N$ is unique, and $\At$ normalizes its image.
\item The preimage of $\{\pm I\}$ is central in $\Gt$.
\end{enumerate}
\end{lemma}

\begin{proof}
(a) We have $x(h(a))=a^{-1}$, $x(h(b))=b^{-1}$ and
$x(h(ab))=(ab)^{-1}$. So $\beta=(b^{-1},a^{-1})_n=(b,a)_n$ by (H1).

(b) Unipotent matrices have $x=1$, so $\beta(u(x),u(y))=(1,1)_n=1$.
Next, $h(t)u(x)$ has lower row $(0,t^{-1})$, so
$x(h(t)u(x))=t^{-1}=x(h(t))$ and $x(u(x))=1$. Thus
$\beta(h(t),u(x))=(1,t^{-1})_n=1$.

(c) Two splittings of $N$ differ by a homomorphism $F\to\mu_n$. Every
element of $F$ is $n$ times another, since $F$ has characteristic
zero. So such a homomorphism $\phi$ satisfies
$\phi(x)=\phi(x/n)^n=1$. Conjugation by an element of $\At$ maps a
splitting of $N$ to a splitting of $N$, because $A$ normalizes $N$. By
uniqueness it preserves the image.

(d) For $\varepsilon=\pm1$ and a lift $z$ of $\varepsilon I$, the
commutator $g\mapsto z\tilde g z^{-1}\tilde g^{-1}$ depends only on
$g=\rho(\tilde g)$. It is a homomorphism $G\to\mu_n$, because $z$ is
central modulo $\mu_n$. The group $\SL(2,F)$ is perfect, so the
homomorphism is trivial.
\end{proof}

Let $K=\SL(2,\Of)$. For $m\ge1$ let $K_m$ be the kernel of reduction
modulo $\varpi^m$.

\paragraph{The splitting of $K$.} Kubota \cite{Ku67} showed that $\Gt$
splits over $K$ when $p\nmid n$. We use the explicit splitting
$s(k)=(k,\kappa(k))$, where
\begin{equation}\label{eq:kappa}
\kappa\begin{pmatrix}a&b\\c&d\end{pmatrix}=
\begin{cases}(c,d)_n,&0<|c|<1,\\ 1,&\text{otherwise,}\end{cases}
\end{equation}
and we prove that it is a splitting. For a unit $u$ let
$\omega(u)\in\mu_n$ be the image of $\bar u^{(q-1)/n}$, where $\bar u$ is
the residue of $u$. Then (H5) reads
$(a,b)_n=\omega\bigl((-1)^{\alpha\beta}a^\beta b^{-\alpha}\bigr)$, with
$\alpha=v(a)$ and $\beta=v(b)$. Note that $\omega$ is multiplicative and
that $\omega(u)=1$ when $u\equiv1\pmod\varpi$.

\begin{lemma}[Kubota's splitting]\label{lem:splitting}
Assume $\gcd(p,n)=1$. For all $k_1,k_2\in K$,
\[
\kappa(k_1k_2)=\kappa(k_1)\kappa(k_2)\beta(k_1,k_2).
\]
Hence $s:K\to\Gt$ is a homomorphism. Moreover $\kappa=1$ on upper
triangular matrices, so $s(b)=(b,1)$ for $b\in B\cap K$.
\end{lemma}

\begin{proof}
\emph{Generators.} The group $K$ is generated by the matrices $u(x)$
with $x\in\Of$ and by $w$. Indeed, let
$k=\left(\begin{smallmatrix}a&b\\c&d\end{smallmatrix}\right)\in K$. If
$c$ is a unit, then
$k=u\bigl((a-1)/c\bigr)\,\bar u(c)\,u\bigl((d-1)/c\bigr)$ with
$\bar u(c)=\left(\begin{smallmatrix}1&0\\c&1\end{smallmatrix}\right)$,
as one checks by multiplying out and using $ad-bc=1$. Also
$\bar u(c)=wu(-c)w^{-1}$ and $w^{-1}=w^3$. If $c$ is not a unit, then
$a$ is a unit, because $ad\equiv1\pmod\varpi$. Then $wk$ has lower row
$(-a,-b)$, so $wk$ is of the first kind, and $k=w^3(wk)$.

\emph{Reduction.} Suppose that $s(gk)=s(g)s(k)$ for every generator
$g$ and every $k\in K$. Then induction on the length of a word gives
$s(g_1\cdots g_rk)=s(g_1)\cdots s(g_r)s(k)$. Taking $k=1$ gives
$s(g_1\cdots g_r)=s(g_1)\cdots s(g_r)$, since $s(1)=(1,1)$ is the
identity. So $s(k_1k_2)=s(k_1)s(k_2)$ for all $k_1,k_2$. In terms of
$\kappa$, the condition $s(gk)=s(g)s(k)$ reads
$\kappa(gk)=\kappa(g)\kappa(k)\beta(g,k)$.

\emph{The generator $u(x)$.} The matrices $u(x)k$ and $k$ have the same
lower row, so $x(u(x)k)=x(k)$ and $\kappa(u(x)k)=\kappa(k)$. Also
$\kappa(u(x))=1$, and $\beta(u(x),k)=(x(k)/1,\,x(k)/x(k))_n=1$.

\emph{The generator $w$.} Here $\kappa(w)=1$, because the lower-left
entry of $w$ is a unit. The lower row of $wk$ is $(-a,-b)$, and
$\beta(w,k)=\bigl(-x(wk),\,x(wk)/x(k)\bigr)_n$, since $x(w)=-1$. We
must show $\kappa(wk)=\kappa(k)\beta(w,k)$. There are four cases.
\begin{enumerate}[label=(\roman*),leftmargin=*]
\item $c$ is a unit and $a$ is $0$ or a unit. Then $x(k)=c$, and
  $x(wk)$ is $-a$ or (if $a=0$) $-b$, a unit since $bc=-1$. All
  arguments of the symbols are units, so both sides equal $1$.
\item $c$ is a unit and $0<|a|<1$. Then $x(k)=c$, $x(wk)=-a$, and $b$ is
  a unit since $bc\equiv-1\pmod\varpi$. Let $\alpha=v(a)$. By (H5),
  $\kappa(wk)=(-a,-b)_n=\omega(-b)^{-\alpha}$, and
  $\beta(w,k)=(a,-a/c)_n=\omega\bigl((-1)^{\alpha}(-c)^{\alpha}\bigr)
  =\omega(c)^{\alpha}$, since $a/(-a/c)=-c$. The two agree because
  $\omega(-b)^{-\alpha}\omega(c)^{-\alpha}=\omega(-bc)^{-\alpha}=1$.
\item $c=0$. Then $a,d$ are units, $x(k)=d$, $x(wk)=-a$, and every
  argument is a unit. Both sides equal $1$.
\item $0<|c|<1$. Then $a,d$ are units with $ad\equiv1\pmod\varpi$,
  $x(k)=c$ and $x(wk)=-a$, so $\kappa(wk)=1$. Let $\gamma=v(c)$. By
  (H5), $\kappa(k)=(c,d)_n=\omega(d)^{-\gamma}$ and
  $\beta(w,k)=(a,-a/c)_n=\omega(a)^{-\gamma}$. Their product is
  $\omega(ad)^{-\gamma}=1$.
\end{enumerate}
This proves the identity. The last assertion holds because $c=0$ for
upper triangular matrices.
\end{proof}

We fix this splitting. We write $K$ also for $s(K)$ when no confusion
can arise. The script \texttt{scripts/kubota\_torus\_check.py} in the
repository also checks Lemma~\ref{lem:splitting} numerically.

\begin{example}\label{ex:noncommute}
Take $F=\Qq_7$ and $n=3$. By Lemma~\ref{lem:cocycle}(a) and
Example~\ref{ex:hilbert},
\[
\hl(7)\hl(3)=(h(21),(3,7)_3),\qquad
\hl(3)\hl(7)=(h(21),(7,3)_3).
\]
The second coordinates are inverse to each other and are nontrivial
cube roots of unity. So $\hl(7)$ and $\hl(3)$ do not commute, although
$h(7)$ and $h(3)$ do.
\end{example}

\subsection{Genuine representations and genuine functions}
\label{subsec:genuine}

A genuine representation and a genuine test function transform by
inverse characters of $\mu_n$, so that their product is invariant and
can be integrated over $\Gt$.

A representation $\pi$ of $\Gt$ is \emph{genuine} if
$\pi(\zeta g)=\zeta\pi(g)$ for $\zeta\in\mu_n$. A function $f$ on
$\Gt$ is \emph{genuine} if
\begin{equation}\label{eq:genuine}
f(\zeta g)=\zeta^{-1}f(g).
\end{equation}
Let $C_c^\infty(\Gt)_\gen$ be the space of genuine, locally constant,
compactly supported functions. For a genuine admissible $\pi$ and
$f\in C_c^\infty(\Gt)_\gen$ the operator
$\pi(f)=\int_{\Gt}f(g)\pi(g)\,dg$ has finite rank. Its integrand is
invariant under $\mu_n$. The character $\Theta_\pi$, where it is a
function, satisfies $\Theta_\pi(\zeta g)=\zeta\Theta_\pi(g)$, and
$\Theta_\pi f$ is invariant under $\mu_n$. The same conventions apply
to representations and functions on $\At$.

\begin{example}\label{ex:eK}
Assume $\gcd(p,2n)=1$ and define the \emph{genuine unit} $e_K$ by
$e_K(\zeta s(k))=\zeta^{-1}$ for $k\in K$, and $e_K=0$ off
$\mu_n s(K)$. It is genuine, and it is invariant under conjugation by
$s(K)$. For a genuine representation $\pi$ and a vector $v$,
\[
\pi(e_K)v=\int_{\mu_n}\int_K\zeta^{-1}\zeta\,\pi(s(k))v\,dk\,d\zeta
=\int_K\pi(s(k))v\,dk .
\]
So $\pi(e_K)$ is the projection onto the $s(K)$-fixed vectors, and
$\tr\pi(e_K)=\dim\pi^{s(K)}$ when $\pi$ is admissible.
\end{example}

\subsection{Haar measures}\label{subsec:haar}

Every identity below depends on Haar measures, so we fix them once
and for all.

Give $\mu_n$ total mass one. On $A\cong F^\times$ use the Haar measure
$dt$ with $\vol(\Of^\times)=1$. On $N\cong F$ use $du$ with
$\vol(\Of)=1$. On $K$ use $dk$ with $\vol(K)=1$. On
$\At=\{(h(t),\zeta)\}$ use the product of $dt$ and $d\zeta$. The
measure on an open subgroup is the \emph{restriction} of the measure
on the group. We never replace it by a pushforward measure without
saying so. For $\gcd(p,2n)=1$ we have $G=ANK$, hence $\Gt=\At N s(K)$,
and we normalize $dg$ on $\Gt$ by
\begin{equation}\label{eq:iwasawa}
\int_{\Gt}\phi(g)\,dg=
\int_{\At}\int_N\int_K\phi(a\,(u,1)\,s(k))\,da\,du\,dk .
\end{equation}
The right side is a Haar measure on $\Gt$ by the standard argument for
Iwasawa decompositions. With this normalization the set $\mu_n s(K)$
has volume one. On $G$ we use the same formula without $\mu_n$. The
quotient $\At\backslash\Gt$ is identified with $A\backslash G$, and in
both cases the quotient measure is $du\,dk$.

\begin{example}\label{ex:volume}
Suppose $p\nmid d$, where $d$ is a positive integer with
$\mu_d\subset F$. The $d$th power map $\Of^\times\to\Of^\times$ has
kernel $\mu_d$ and image $(\Of^\times)^d$. Hence $(\Of^\times)^d$ has
index $d$ in $\Of^\times$, and its restricted volume is $1/d$. For
instance, $(\Zz_7^\times)^3$ has index $3$ in $\Zz_7^\times$ and
volume $1/3$. This factor appears in Lemma~\ref{lem:norm}.
\end{example}

\subsection{Orbital integrals on \texorpdfstring{$\SL(2,F)$}{SL(2,F)}}\label{subsec:orbital-linear}

For the linear group, split orbital integrals are constant terms along
$N$, and the Weyl integration formula expresses integrals over $G$
through them. Section~\ref{sec:orbital} generalizes both to the
cover.

Let $a=h(t)$ with $t^2\ne1$. The \emph{Weyl discriminant} and the
\emph{modulus character} are
\begin{align}
D(a)&=\det\bigl(1-\Ad(a)\bigr)\big|_{\mathfrak g/\mathfrak a}
     =(1-t^2)(1-t^{-2})=-(t-t^{-1})^2,\label{eq:D}\\
\delta_B(a)&=|t|^2,\qquad |D(a)|^{1/2}=|t-t^{-1}|.\label{eq:delta}
\end{align}
Here $\mathfrak g$ and $\mathfrak a$ are the Lie algebras of $G$ and
$A$. Both quantities pull back to $\Gt$ through $\rho$. For
$f\in C_c^\infty(G)$ put $f^K(g)=\int_Kf(k^{-1}gk)\,dk$. The
\emph{normalized split orbital integral} of $f$ is
\begin{equation}\label{eq:FA}
F_f^A(t)=|D(h(t))|^{1/2}\int_{A\backslash G}f(x^{-1}h(t)x)\,d\dot x,
\qquad t^2\ne1 .
\end{equation}
The change of variables in the proof of Lemma~\ref{lem:abel} below,
with $n=1$, gives
\begin{equation}\label{eq:abel-linear}
F_f^A(t)=|t|\int_Nf^K(h(t)u)\,du .
\end{equation}
The right side is defined for all $t$, is locally constant and is
compactly supported. It is the locally constant extension of $F_f^A$.
It is invariant under $t\mapsto t^{-1}$.

The \emph{Weyl integration formula} \cite{Si79} states that for
integrable $\phi$ on $G$,
\begin{equation}\label{eq:weyl}
\int_G\phi(g)\,dg=\frac12\int_A|D(a)|\int_{A\backslash G}
\phi(x^{-1}ax)\,d\dot x\,da+E(\phi),
\end{equation}
where $E(\phi)$ is a finite sum of similar integrals over the elliptic
maximal tori. It depends only on the restriction of $\phi$ to the set
$G_{\mathrm{ell}}$ of regular elliptic elements.

\begin{example}\label{ex:FA-unit}
Let $f=1_K$. Then $f^K=1_K$. For a unit $t$, $h(t)u(x)\in K$ if and
only if $x\in\Of$. For a nonunit $t$, $h(t)u(x)\notin K$. By
\eqref{eq:abel-linear},
\[
F_{1_K}^A(t)=1_{\Of^\times}(t).
\]
\end{example}

\subsection{Principal series of \texorpdfstring{$\SL(2,F)$}{SL(2,F)}}\label{subsec:ps-linear}

The trace of a linear principal series is the integral of its
inducing character against the split orbital integral.
Theorem~\ref{thm:transfer} compares this with the cover.

Let $\eta$ be a smooth character of $A\cong F^\times$. The normalized
principal series $\Ind_B^G(\eta\delta_B^{1/2})$ acts by right
translation on locally constant $\phi:G\to\Cc$ with
$\phi(bg)=\delta_B(b)^{1/2}\eta(b)\phi(g)$. Its character is
\begin{equation}\label{eq:char-linear}
\Theta(h(t))=\frac{\eta(t)+\eta(t)^{-1}}{|t-t^{-1}|}\quad(t^2\ne1),
\qquad \Theta=0\text{ on }G_{\mathrm{ell}},
\end{equation}
and its trace is
\begin{equation}\label{eq:trace-linear}
\tr\Ind_B^G(\eta\delta_B^{1/2})(f)=\int_A\eta(t)F_f^A(t)\,dt .
\end{equation}
Both formulas are the case $n=1$ of Lemma~\ref{lem:induced-trace},
Corollary~\ref{cor:trace-integral} and Theorem~\ref{thm:character}
below, so the paper proves them.

\begin{example}\label{ex:trace-linear}
Let $\eta$ be unramified and $f=1_K$. By Example~\ref{ex:FA-unit} and
\eqref{eq:trace-linear}, the trace is $\int_{\Of^\times}\eta\,dt=1$.
This is the dimension of the space of $K$-fixed vectors.
\end{example}

%======================================================================
\section{Notation}\label{sec:notation}

The table collects the principal notation. Definitions are in
Sections~\ref{sec:background} and~\ref{sec:torus}. For example, when
$n=6$ one has $d=3$, $[\At:\ZA]=9$ and $\dim\chi_\lambda=3$.

\begin{center}
\begingroup
\small
\setlength{\tabcolsep}{4pt}
\renewcommand{\arraystretch}{1.15}
\begin{tabularx}{\linewidth}{@{}>{\raggedright\arraybackslash}p{0.27\linewidth}>{\raggedright\arraybackslash}X@{}}
\toprule
Symbol & Meaning\\
\midrule
$F,\Of,\varpi,q,|\cdot|,v$ & Local field, integers, uniformizer,
residue field order, absolute value, valuation.\\
$n$, $\mu_n$, $(\ ,\ )_n$ & Degree of the cover, $n$th roots of unity,
$n$th Hilbert symbol (Section~\ref{subsec:hilbert}).\\
$d=n/\gcd(n,2)$ & Dimension of a genuine irreducible representation
of $\At$.\\
$G$, $\Gt=\overline{SL(2,F)}$, $\rho$ & $\SL(2,F)$, Kubota's $n$-fold
cover, the projection.\\
$h(t),u(x),w$ & Diagonal, unipotent and Weyl elements of $G$.\\
$A,N,B,W$ & Diagonal torus, upper unipotent group, Borel subgroup,
Weyl group.\\
$\At,\Bt$, $\hl(t),\wt$ & Preimages of $A$ and $B$. The lifts
$(h(t),1)$ and $(w,1)$.\\
$K$, $K_m$, $s$, $\kappa$ & $\SL(2,\Of)$, its principal congruence
subgroups, the splitting \eqref{eq:kappa}.\\
$\ZA$, $\ZA^0$ & Center of $\At$, equal to
$\rho^{-1}(h((F^\times)^d))$. The set
$\{(h(u),1):u\in(F^\times)^d\}$.\\
$C$, $\Ct$ & $h(\varpi^{d\Zz}\Of^\times)$ and its preimage, a maximal
abelian subgroup of $\At$.\\
$\lambda$, $\tilde\lambda$, $\chi_\lambda$ & Genuine character of
$\ZA$, an extension to $\Ct$, the induced representation of $\At$.\\
$\lambda^w$, $\chi^{\wt}$ & Weyl conjugates
$z\mapsto\lambda(\wt^{-1}z\wt)$ and $a\mapsto\chi(\wt^{-1}a\wt)$.\\
$D,\delta_B$ & Weyl discriminant and modulus character.\\
$f^K$, $F_f^A$ & $K$-average and linear split orbital integral.\\
$\Psi_f$, $\Phi_f$, $\mathcal O_f^{\ZA}$ & Constant term on $\At$ and
normalized split orbital integrals on $\ZA$ and $\At$
(Section~\ref{sec:orbital}).\\
$e_K$ & Genuine unit (Example~\ref{ex:eK}).\\
$\PI(\chi)$, $\Theta_\Pi$ & Normalized genuine principal series and
its character.\\
$\nu$, $\eta_\lambda$ & Power map $h(t)\mapsto(h(t^d),1)$ and the
character $t\mapsto\lambda(\nu(h(t)))$ of $A$.\\
\bottomrule
\end{tabularx}
\endgroup
\end{center}

%======================================================================
\section{The covering torus}\label{sec:torus}

The commutator of two lifts of diagonal elements is a power of the
Hilbert symbol. Its radical is the center $\ZA$, and a Stone--von
Neumann argument parametrizes the genuine irreducible representations
of $\At$ by genuine characters of $\ZA$.

\subsection{The commutator, the center and Weyl conjugation}

We compute the commutator pairing on $\At$, its radical, and the
action of $\wt$.

\begin{lemma}\label{lem:commutator}
For $a,b\in F^\times$,
\begin{equation}\label{eq:commutator}
\hl(a)\hl(b)\hl(a)^{-1}\hl(b)^{-1}=(a,b)_n^{-2}.
\end{equation}
\end{lemma}

\begin{proof}
By Lemma~\ref{lem:cocycle}(a), $\hl(a)\hl(b)=(h(ab),(b,a)_n)$ and
$\hl(b)\hl(a)=(h(ab),(a,b)_n)$. Hence
$\hl(a)\hl(b)=(b,a)_n(a,b)_n^{-1}\,\hl(b)\hl(a)=(a,b)_n^{-2}\,
\hl(b)\hl(a)$ by (H2).
\end{proof}

Define
\begin{equation}\label{eq:Z}
\ZA=\rho^{-1}\bigl(h((F^\times)^d)\bigr),\qquad d=\frac{n}{\gcd(n,2)} .
\end{equation}
This is a set of the form $\mu_n\cdot\ZA^0$ with
$\ZA^0=\{(h(u),1):u\in(F^\times)^d\}$. It is not the set of $d$th
powers of elements of $\At$.

\begin{lemma}\label{lem:radical}
The group $\ZA$ is the center of $\At$. If $\gcd(p,n)=1$, then
\begin{equation}\label{eq:indices}
[\At:\ZA]=d^2,
\end{equation}
and the commutator induces a nondegenerate alternating pairing on the
finite group $\At/\ZA$.
\end{lemma}

\begin{proof}
Since $\mu_n$ is central, $\hl(a)$ is central in $\At$ exactly when
$(a,b)_n^2=1$ for all $b$, by Lemma~\ref{lem:commutator}. By (H1) this
says $(a^2,b)_n=1$ for all $b$. By (H4) it says $a^2\in(F^\times)^n$.
We show that this holds exactly when $a\in(F^\times)^d$.

If $n$ is odd, then $d=n$. Choose integers $k,m$ with $2k+nm=1$. Then
$a=(a^2)^k(a^n)^m$, so $a^2\in(F^\times)^n$ implies
$a\in(F^\times)^n$. The converse is clear.

If $n$ is even, then $n=2d$. If $a^2=c^{2d}$, then $(a/c^d)^2=1$, so
$a=\pm c^d$. A primitive $n$th root of unity $\zeta$ lies in $F$, and
$\zeta^d=-1$. So $-1\in(F^\times)^d$ and $a\in(F^\times)^d$.
Conversely, $a=c^d$ gives $a^2=c^n$.

Now let $p\nmid n$. Since $\mu_n\subset F$ and $p\nmid n$, the group
$\mu_n$ injects into the residue field, so $n\mid q-1$ and $d\mid q-1$.
We have $F^\times\cong\varpi^\Zz\times\mu_{q-1}\times(1+\varpi\Of)$.
The $d$th power map is bijective on the pro-$p$ group $1+\varpi\Of$.
Hence $F^\times/(F^\times)^d\cong\Zz/d\times\Zz/d$, which gives
\eqref{eq:indices}. The pairing $(a,b)\mapsto(a,b)_n^{-2}$ is
alternating because $(a,a)_n^2=(a,-1)_n^2=1$. It is nondegenerate on
$\At/\ZA$ by the first part of the proof.
\end{proof}

\begin{remark}\label{rem:centralizer}
The centralizer of all of $\At$ is $\ZA$. The centralizer in $\Gt$ of
a single regular element $a\in\At$ can be larger than $\ZA$ and
smaller than $\At$. So one may not integrate over $\At\backslash\Gt$
in an orbital integral unless $a$ commutes with all of $\At$, that is,
unless $a\in\ZA$.
\end{remark}

\begin{lemma}\label{lem:weyl-torus}
For $t\in F^\times$:
\begin{enumerate}[label=(\alph*)]
\item $\wt^{-1}=(w^{-1},1)$.
\item $\wt^{-1}\hl(t)\wt=\hl(t^{-1})$.
\item $\hl(t)^{-1}=(t,-1)_n\,\hl(t^{-1})$.
\item $\wt^2$ lies in the preimage of $-I$, so it is central in $\Gt$.
\end{enumerate}
\end{lemma}

\begin{proof}
(a) We have $x(w)=-1$, $x(w^{-1})=1$ and $x(I)=1$, so
$\beta(w,w^{-1})=(-1,1)_n=1$.

(b) The matrix $w^{-1}h(t)$ has lower row $(t,0)$. So
$x(w^{-1}h(t))=t$, while $x(h(t))=t^{-1}$. Thus
$\beta(w^{-1},h(t))=(t,t^2)_n=(t,t)_n^2=(t,-1)_n^2=1$. Next,
$w^{-1}h(t)w=h(t^{-1})$ and $x(h(t^{-1}))=t$. Thus
\[
\beta(w^{-1}h(t),w)=\Bigl(\frac tt,\ \frac t{-1}\Bigr)_n=(1,-t)_n=1 .
\]
Multiplying out,
$(w^{-1},1)(h(t),1)(w,1)=(h(t^{-1}),1)$.

(c) By Lemma~\ref{lem:cocycle}(a), $\hl(t)\hl(t^{-1})=(I,(t^{-1},t)_n)$
and $(t^{-1},t)_n=(t,t)_n^{-1}=(t,-1)_n^{-1}=(t,-1)_n$. Hence
$\hl(t)^{-1}=(t,-1)_n\hl(t^{-1})$.

(d) $w^2=-I$, and Lemma~\ref{lem:cocycle}(d) applies.
\end{proof}

By (b), Weyl conjugation and inversion agree on the matrix part of
$\hl(t)$. They differ on $\mu_n$: $\wt^{-1}\zeta\hl(t)\wt=\zeta\hl(t^{-1})$,
whereas $(\zeta\hl(t))^{-1}=\zeta^{-1}(t,-1)_n\hl(t^{-1})$.

\begin{example}\label{ex:inverse}
Take $F=\Qq_5$, $n=4$ and $t=5$. By Example~\ref{ex:hilbert},
$(5,-1)_4=-1$. So $\hl(5)^{-1}=-\hl(1/5)$, while
$\wt^{-1}\hl(5)\wt=\hl(1/5)$. Even on the lifts $\hl(t)$, inversion
and Weyl conjugation differ.
\end{example}

\subsection{The Stone--von Neumann parametrization}\label{subsec:svn}

Because the commutator pairing on $\At/\ZA$ is nondegenerate, a maximal
abelian subgroup has index $d$, and every genuine irreducible
representation is induced from one of its characters. This is the
rank-one analogue of the torus theory of \cite{Mc12}. Since that work
assumes $\mu_{2n}\subset F$, we give the argument in full.
Assume $\gcd(p,n)=1$ and put
\[
C=h(\varpi^{d\Zz}\Of^\times),\qquad \Ct=\rho^{-1}(C).
\]
Two elements $h(\varpi^{dk}u)$ and $h(\varpi^{dl}u')$ of $C$ have
commutator $(\varpi^{dk}u,\varpi^{dl}u')_n^{-2}$. By (H1) and (H5) this
equals $(\varpi,u')_n^{-2dk}(u,\varpi)_n^{-2dl}(\varpi,\varpi)_n^{-2d^2kl}$.
Each factor is an $n$th root of unity raised to a multiple of $2d$,
which is a multiple of $n$. So $\Ct$ is abelian. It contains $\ZA$, and
$[\At:\Ct]=[F^\times:\varpi^{d\Zz}\Of^\times]=d$. Since $d^2=[\At:\ZA]$
and the pairing on $\At/\ZA$ is nondegenerate, $\Ct$ is a maximal
abelian subgroup.

\begin{theorem}[Genuine representations of $\At$]\label{thm:torus}
Assume $\gcd(p,n)=1$. For each smooth genuine character
$\lambda:\ZA\to\Cc^\times$ choose an extension $\tilde\lambda$ to
$\Ct$. Then
\begin{equation}\label{eq:stone}
\chi_\lambda=\Ind_{\Ct}^{\At}\tilde\lambda
\end{equation}
is irreducible of dimension $d$. Its isomorphism class does not depend
on the extension. The map $\lambda\mapsto\chi_\lambda$ is a bijection
from genuine smooth characters of $\ZA$ to genuine smooth irreducible
representations of $\At$. The trace is
\begin{equation}\label{eq:torus-trace}
\tr\chi_\lambda(a)=
\begin{cases}
d\,\lambda(a),&a\in\ZA,\\
0,&a\in\At\setminus\ZA.
\end{cases}
\end{equation}
\end{theorem}

\begin{proof}
An extension exists because $\Ct/\ZA$ is a finite abelian group and
$\Cc^\times$ is divisible. Every extension is smooth, because $\ZA$ is
open in $\Ct$ and $\lambda$ is smooth. For $b\in\At$ the conjugate extension is
$\tilde\lambda^b(c)=\tilde\lambda(b^{-1}cb)=\langle b,c\rangle\,
\tilde\lambda(c)$, where $\langle b,c\rangle=b^{-1}cbc^{-1}\in\mu_n$
is a commutator. The map $b\mapsto\langle b,\cdot\rangle$ embeds
$\At/\Ct$ into the character group of $\Ct/\ZA$, by maximality of
$\Ct$. Both groups have order $d$. So the conjugates of
$\tilde\lambda$ are exactly the $d$ extensions of $\lambda$, and the
stabilizer of $\tilde\lambda$ is $\Ct$. Mackey's criterion shows that
$\chi_\lambda$ is irreducible of dimension $[\At:\Ct]=d$, and that its
class does not depend on the extension.

Conversely, let $\sigma$ be a genuine smooth irreducible
representation of $\At$. By Schur's lemma $\ZA$ acts by a genuine
character $\lambda$. Since $\Ct/\ZA$ is finite and abelian, the
restriction of $\sigma$ to $\Ct$ contains a character $\tilde\lambda$
extending $\lambda$. Frobenius reciprocity gives a nonzero map
$\chi_\lambda\to\sigma$, which is an isomorphism because both are
irreducible.

For $a\in\ZA$, Schur's lemma gives $\chi_\lambda(a)=\lambda(a)$ times
the identity, so the trace is $d\lambda(a)$. If $a\notin\ZA$, choose
$b\in\At$ with $b^{-1}ab=\zeta a$ and $\zeta\ne1$. Then
$\tr\chi_\lambda(a)=\tr\chi_\lambda(b^{-1}ab)=\zeta\tr\chi_\lambda(a)$,
so the trace vanishes.
\end{proof}

\begin{example}[The double and triple covers]\label{ex:two-three}
For $n=2$ we have $d=1$. The group $\At$ is abelian, $\ZA=\At$, and
the genuine irreducible representations of $\At$ are characters. Still,
$t\mapsto\hl(t)$ is not a homomorphism, since
$\beta(h(a),h(b))=(b,a)_2$ need not be $1$.

For $n=3$ we have $d=3$ and $[\At:\ZA]=9$. Each genuine irreducible
representation of $\At$ is three-dimensional and is induced from a
character of a subgroup of index three. Its trace is supported on
$\ZA$ and equals three times the central character there. By
Example~\ref{ex:noncommute}, with $F=\Qq_7$, the element $\hl(7)$ is not
in $\ZA$, and its trace is zero.
\end{example}

\begin{example}[Degrees four and six]\label{ex:four-six}
For $n=4$ we have $d=2$ and $[\At:\ZA]=4$. The genuine irreducible
representations of $\At$ are two-dimensional. The map
$t\mapsto(h(t^2),1)$ is a homomorphism, because
$\beta(h(a^2),h(b^2))=(b,a)_4^4=1$. So Theorem~\ref{thm:transfer}
applies in degree four.

For $n=6$ we have $d=3$ and $[\At:\ZA]=9$. Now
$\beta(h(a^3),h(b^3))=(b,a)_6^9=(b,a)_6^3$, which need not be $1$.
Theorem~\ref{thm:torus} still holds, but the map
$t\mapsto(h(t^3),1)$ is not a homomorphism. This is why
Theorem~\ref{thm:transfer} excludes $n=6$.
\end{example}

%======================================================================
\section{Split orbital integrals and their image}\label{sec:orbital}

Only elements of the center $\ZA$ have nonzero normalized orbital
integrals. On $\ZA$ these integrals fill out exactly the Weyl-invariant
genuine test functions. That image theorem is what makes the comparison
with $\SL(2,F)$ in Section~\ref{sec:transfer} possible.

\subsection{The correct quotient}

Because $\At$ is not abelian, the group divided out in an orbital
integral must centralize the element (Remark~\ref{rem:centralizer}).

Since $\ZA$ centralizes $\At$, the integral
\begin{equation}\label{eq:OZ}
\mathcal O_f^{\ZA}(a)=|D(a)|^{1/2}\int_{\ZA\backslash\Gt}
f(x^{-1}ax)\,d\dot x,\qquad a\in\At_\reg,
\end{equation}
is well defined for $f\in C_c^\infty(\Gt)_\gen$. Here
$\At_\reg=\rho^{-1}(A_\reg)$ and $A_\reg=\{h(t):t^2\ne1\}$. For
$z\in(\ZA)_\reg=\ZA\cap\At_\reg$ the centralizer of $z$ in $\Gt$ is
$\At$, so we also set
\begin{equation}\label{eq:Phi}
\Phi_f(z)=|D(z)|^{1/2}\int_{\At\backslash\Gt}f(x^{-1}zx)\,d\dot x .
\end{equation}
The measure on $\ZA$ is the restriction of the measure on $\At$, so
$\ZA\backslash\At$ has volume $[\At:\ZA]$. Hence
\begin{equation}\label{eq:index-phi}
\mathcal O_f^{\ZA}(z)=d^2\,\Phi_f(z)\qquad(z\in(\ZA)_\reg)
\end{equation}
when $\gcd(p,n)=1$. The index $d^2$ should not be confused with $d$,
the dimension of $\chi_\lambda$.

\paragraph{Three quotients.} Three subgroups of $\At$ centralize an
element $z\in(\ZA)_\reg$: the center $\ZA$, the maximal abelian subgroup
$\Ct$, and $\At$ itself. Since $\ZA\subset\Ct\subset\At$ with
$[\Ct:\ZA]=[\At:\Ct]=d$, the three orbital integrals of $f$ at $z$
differ only by these indices:
\begin{equation}\label{eq:three-quotients}
\begin{array}{lll}
\text{quotient} & \text{index over }\At\backslash\Gt & \text{normalized integral at }z\\[2pt]
\ZA\backslash\Gt & d^2 & \mathcal O_f^{\ZA}(z)=d^2\,\Phi_f(z)\\
\Ct\backslash\Gt & d & \mathcal O_f^{\Ct}(z)=d\,\Phi_f(z)\\
\At\backslash\Gt & 1 & \Phi_f(z)
\end{array}
\end{equation}
Here $\mathcal O_f^{\Ct}$ is defined like \eqref{eq:OZ} with $\Ct$ in
place of $\ZA$. For $a\in\At_\reg\setminus\ZA$ the $\ZA$-quotient
integral is still defined, and the $\Ct$-quotient integral is defined
when $a\in\Ct$. Both vanish there, by Lemma~\ref{lem:vanishing} and the
same proof, since $\Ct$ is normal in $\At$.

\begin{lemma}[Vanishing]\label{lem:vanishing}
For $f\in C_c^\infty(\Gt)_\gen$ and $a\in\At_\reg\setminus\ZA$,
$\mathcal O_f^{\ZA}(a)=0$.
\end{lemma}

\begin{proof}
By Lemma~\ref{lem:radical} there is $b\in\At$ with $b^{-1}ab=\zeta a$
and $\zeta\ne1$. Left translation $x\mapsto bx$ is well defined on
$\ZA\backslash\Gt$, because $b$ centralizes $\ZA$. It preserves the
quotient measure, because $\At$ and $\Gt$ are unimodular. After the
substitution the integrand is
$f(x^{-1}b^{-1}abx)=f(\zeta x^{-1}ax)=\zeta^{-1}f(x^{-1}ax)$. Hence
$\mathcal O_f^{\ZA}(a)=\zeta^{-1}\mathcal O_f^{\ZA}(a)$, and the
integral is zero.
\end{proof}

\begin{example}\label{ex:vanishing}
Take $F=\Qq_7$, $n=3$ and $a=\hl(7)$, which is regular. By
Example~\ref{ex:noncommute}, $b=\hl(3)$ satisfies $b^{-1}ab=\zeta a$
with $\zeta$ a nontrivial cube root of unity. So
$\mathcal O_f^{\ZA}(\hl(7))=0$ for every genuine $f$.
\end{example}

\subsection{The constant term and the image theorem}

Writing $\Phi_f$ as a constant term along $N$ extends it across the
singular set and leads to the image theorem. Assume $\gcd(p,2n)=1$ and use the splitting $s$ of $K$. For
$f\in C_c^\infty(\Gt)_\gen$ put $f^K(g)=\int_Kf(s(k)^{-1}gs(k))\,dk$,
and for $a\in\At$ put
\begin{equation}\label{eq:Psi}
\Psi_f(a)=\delta_B(a)^{1/2}\int_Nf^K(a\,(u,1))\,du .
\end{equation}

\begin{lemma}\label{lem:abel}
Assume $\gcd(p,2n)=1$. Then $\Psi_f$ is locally constant and compactly
supported on $\At$, and $\Psi_f(\zeta a)=\zeta^{-1}\Psi_f(a)$. For
$z\in(\ZA)_\reg$,
\begin{equation}\label{eq:Abel}
\Phi_f(z)=\Psi_f(z)=\delta_B(z)^{1/2}\int_Nf^K(z\,(u,1))\,du .
\end{equation}
\end{lemma}

\begin{proof}
The first assertions follow because $f^K$ is genuine, locally constant
and compactly supported. Let $z\in(\ZA)_\reg$ with $\rho(z)=h(t)$. By
\eqref{eq:iwasawa} the quotient measure on $\At\backslash\Gt$ is
$du\,dk$, so
\[
\Phi_f(z)=|D(z)|^{1/2}\int_Nf^K\bigl((u,1)^{-1}z(u,1)\bigr)du .
\]
Write $u=u(x)$. By Lemma~\ref{lem:cocycle}(c), $z^{-1}(u(x),1)z$ is
the canonical lift of $h(t)^{-1}u(x)h(t)=u(t^{-2}x)$. Hence
$(u(x),1)^{-1}z(u(x),1)=z\,(u((1-t^{-2})x),1)$. Substitute
$y=(1-t^{-2})x$, so that $dx=|1-t^{-2}|^{-1}dy$. Since
$|D(z)|^{1/2}=|t|\,|1-t^{-2}|$ and $\delta_B(z)^{1/2}=|t|$, we obtain
\eqref{eq:Abel}.
\end{proof}

We use $\Psi_f|_{\ZA}$ as the locally constant extension of $\Phi_f$
to all of $\ZA$, and we keep the name $\Phi_f$. On $(\ZA)_\reg$,
$\Phi_f$ is invariant under $z\mapsto\wt^{-1}z\wt$. Indeed $\wt$
normalizes $\At$, so $x\mapsto\wt x$ preserves the measure on
$\At\backslash\Gt$, and $|D|$ is invariant. By local constancy the
extension is also invariant. Let $C_c^\infty(\ZA)_\gen^W$ be the space
of locally constant, compactly supported functions $\phi$ on $\ZA$ with
$\phi(\zeta z)=\zeta^{-1}\phi(z)$ and $\phi(\wt^{-1}z\wt)=\phi(z)$. By
Lemma~\ref{lem:weyl-torus}(b) the second condition reads
$\phi(\zeta\hl(u^{-1}))=\phi(\zeta\hl(u))$.

\begin{theorem}[Image of the split orbital integral]\label{thm:orbital}
Assume $\gcd(p,2n)=1$. The map
\begin{equation}\label{eq:orb-surj}
C_c^\infty(\Gt)_\gen\longrightarrow C_c^\infty(\ZA)_\gen^W,\qquad
f\longmapsto\Phi_f,
\end{equation}
is well defined and surjective.
\end{theorem}

\begin{proof}
Lemma~\ref{lem:abel} and the remarks after it show that the map lands
in $C_c^\infty(\ZA)_\gen^W$. We prove surjectivity in two steps.

\emph{Step 1: functions supported on the regular set.} Consider
\[
\mu:(\At\backslash\Gt)\times(\ZA)_\reg\to\Gt,\qquad
(\At x,z)\mapsto x^{-1}zx .
\]
It is well defined because $z$ is central in $\At$. Its composition
with $\rho$ is the map $(Ax,h(t))\mapsto x^{-1}h(t)x$. At a regular
point the differential of that map is bijective, since
$\Ad(h(t))-1$ is invertible on $\mathfrak g/\mathfrak a$. By the
$p$-adic inverse function theorem it is a local homeomorphism onto an
open set. Since $\rho$ is a local homeomorphism, so is $\mu$.

We describe the fibers of $\mu$. Suppose $x^{-1}zx=y^{-1}z'y$ and put
$c=yx^{-1}$. Then $c$ conjugates $z$ to $z'$. Both are regular, so
$\rho(c)$ normalizes $A$ and $c\in\At\cup\wt\At$. If $c\in\At$, then
$z'=z$ and $\At y=\At x$. If $c=\wt b$ with $b\in\At$, then
$z'=\wt z\wt^{-1}$ and $\At y=\At\wt x$. So each fiber is
$\{(\At x,z),(\At\wt x,\wt z\wt^{-1})\}$.

Let $\alpha$ be the characteristic function of $\At s(K)$, viewed on
$\At\backslash\Gt$. By \eqref{eq:iwasawa}, $\int\alpha=1$. Let
$\phi\in C_c^\infty(\ZA)_\gen^W$ vanish near the singular set. Define
$f$ on the image of $\mu$ by
\[
f(x^{-1}zx)=\phi(z)\bigl(\alpha(\At x)+\alpha(\At\wt x)\bigr),
\]
and $f=0$ elsewhere. The value at the second point of the fiber is
$\phi(\wt z\wt^{-1})(\alpha(\At\wt x)+\alpha(\At\wt^2x))$. It is the
same, because $\phi$ is Weyl invariant and $\wt^2\in\At$. So $f$ is well
defined. It is genuine because $\phi$ is.

We check that $f\in C_c^\infty(\Gt)$. Let
$S=\operatorname{supp}\alpha\cup\wt^{-1}\operatorname{supp}\alpha$, a
compact open subset of $\At\backslash\Gt$, and let
$\Sigma=\operatorname{supp}\phi$, a compact open subset of $(\ZA)_\reg$.
The function $f\circ\mu$ is locally constant and vanishes off
$S\times\Sigma$. Since $\mu$ is a local homeomorphism, $f$ is locally
constant on the open set $U=\mu((\At\backslash\Gt)\times(\ZA)_\reg)$.
The support of $f$ lies in $\mu(S\times\Sigma)$, which is compact
because $\mu$ is continuous. A compact subset of the open set $U$ has a
neighborhood on which $f$ is defined by the formula, and $f$ vanishes
on the open complement of $\mu(S\times\Sigma)$. These two open sets
cover $\Gt$, and $f$ is locally constant on each. So $f$ is locally
constant on $\Gt$, and it has compact support. Finally, the
left translation $x\mapsto\wt x$ preserves the measure on
$\At\backslash\Gt$, so
\[
\Phi_f(z)=|D(z)|^{1/2}\phi(z)\int\bigl(\alpha(\At x)+\alpha(\At\wt x)
\bigr)d\dot x=2|D(z)|^{1/2}\phi(z).
\]
Applying this to $\phi/(2|D|^{1/2})$, which lies in the same space,
gives a preimage of $\phi$.

\emph{Step 2: the singular set.} The singular elements of $\ZA$ lie over
$\pm I$. For $\varepsilon=\pm1$ let $z_\varepsilon=(\varepsilon I,1)$,
which is central by Lemma~\ref{lem:cocycle}(d), and lies in $\ZA$
because $-1\in(F^\times)^d$ (proof of Lemma~\ref{lem:radical}). For
$m\ge1$ let
$f_{\varepsilon,m}(\zeta z_\varepsilon s(k))=\zeta^{-1}$ for
$k\in K_m$, and $f_{\varepsilon,m}=0$ off $\mu_nz_\varepsilon s(K_m)$.
This function is genuine, locally constant and compactly supported. It
is invariant under conjugation by $s(K)$, because $z_\varepsilon$ is
central and $K_m$ is normal in $K$. So $f_{\varepsilon,m}^K=
f_{\varepsilon,m}$.

We compute $\Phi_{f_{\varepsilon,m}}$ on all of $\ZA$. Let $z\in\ZA$ and
$x\in F$. If $z(u(x),1)\in\mu_nz_\varepsilon s(K_m)$, then
$\rho(z)u(x)\in\varepsilon K_m$, which forces
$\rho(z)\in\varepsilon h(1+\varpi^m\Of)$ and $x\in\varpi^m\Of$. So
$\Phi_{f_{\varepsilon,m}}$ vanishes off the set
$U_{\varepsilon,m}=\mu_nz_\varepsilon\{(h(t),1):t\in1+\varpi^m\Of\}$.
Every element of $U_{\varepsilon,m}$ has the form
$z=\zeta z_\varepsilon(h(t),1)$. By Lemma~\ref{lem:cocycle}(b) and
Lemma~\ref{lem:splitting}, $z(u(x),1)=\zeta z_\varepsilon s(h(t)u(x))$,
and this lies in $\mu_nz_\varepsilon s(K_m)$ exactly when
$x\in\varpi^m\Of$. By \eqref{eq:Abel},
\[
\Phi_{f_{\varepsilon,m}}=q^{-m}\,\mathbf 1^{\gen}_{U_{\varepsilon,m}}
\quad\text{on all of }\ZA,
\]
where $\mathbf 1^{\gen}_{U}(\zeta z')=\zeta^{-1}$ for $\zeta z'\in U$ with
$z'\in z_\varepsilon\{(h(t),1)\}$, and $\mathbf 1^{\gen}_U=0$ off $U$. In
the tame case $1+\varpi^m\Of\subset(F^\times)^d$, so $U_{\varepsilon,m}$
is an open neighborhood of $z_\varepsilon$ in $\ZA$. It is Weyl stable,
by Lemma~\ref{lem:weyl-torus}(b). Since $p$ is odd, $-1\notin
1+\varpi\Of$, so $U_{1,m}$ and $U_{-1,m}$ are disjoint for every
$m\ge1$.

Now let $\phi\in C_c^\infty(\ZA)_\gen^W$ be arbitrary. Choose $m$ so
large that $\phi(\zeta z_\varepsilon(h(t),1))=\zeta^{-1}
\phi(z_\varepsilon)$ for $t\in1+\varpi^m\Of$ and $\varepsilon=\pm1$.
Then
\[
\phi-\sum_{\varepsilon=\pm1}q^m\phi(z_\varepsilon)\,
\Phi_{f_{\varepsilon,m}}
\]
vanishes near the singular set and lies in $C_c^\infty(\ZA)_\gen^W$.
Step~1 gives a preimage of it, and so $\phi$ has a preimage.
\end{proof}

\begin{remark}[Comparison with the 1999 and 2001 normalizations]
\label{rem:image-jo01}
The 1999 manuscript writes the transform as $F_f^{\overline A}$ on the
preimage of $A^{n_0}=h((F^\times)^d)$, which is $\ZA$, with the quotient
$\overline A\backslash\overline G$. With the measures of
Section~\ref{subsec:haar} it is our $\Phi_f$. If an orbital integral is
formed instead with the quotient by $\ZA$ or by the maximal abelian
subgroup $\Ct$, then \eqref{eq:three-quotients} gives the conversion:
multiply by $d^2$ or by $d$. In each normalization the image is the
same space $C_c^\infty(\ZA)_\gen^W$. The image theorem of \cite{Jo01}
uses yet another notation, and its quotient and measures must be
matched by the same table before the two statements are identified.
\end{remark}

\begin{example}\label{ex:Phi-unit}
Assume $\gcd(p,2n)=1$ and let $e_K$ be the genuine unit of
Example~\ref{ex:eK}. Every element of $\ZA$ is $\zeta(h(u),1)$ with
$u\in(F^\times)^d$. For $u\in\Of^\times$ and $x\in\Of$,
$(h(u),1)(u(x),1)=s(h(u)u(x))$ by Lemma~\ref{lem:cocycle}(b). For
$u\notin\Of^\times$ the matrix $h(u)u(x)$ is not in $K$. Since
$e_K^K=e_K$, \eqref{eq:Abel} gives
\[
\Phi_{e_K}(\zeta(h(u),1))=\zeta^{-1}1_{\Of^\times}(u)
\qquad(u\in(F^\times)^d).
\]
Note that $(F^\times)^d\cap\Of^\times=(\Of^\times)^d$.
\end{example}

\begin{remark}[Satake and the image theorem]\label{rem:satake}
Formula~\eqref{eq:Abel} is a constant term of the conjugation average
of a test function. It should not be confused with the spherical
Satake isomorphism of convolution algebras \cite{Mc12}. Surjectivity of
\eqref{eq:orb-surj} on all test functions does not by itself give a
statement about spherical functions.
\end{remark}

%======================================================================
\section{Principal series and known unitary representations}
\label{sec:unitary}

Genuine principal series are induced from the $d$-dimensional
representations $\chi_\lambda$ of $\At$. We compute their Jacquet
modules and prove unitarity, an irreducibility criterion and the
structure of the spherical vectors. We then survey the literature on
reducibility, complementary series and supercuspidal representations.
Statements proved elsewhere are in remarks numbered ``Remark 2026.$n$''
with their status. The examples treat the threefold cover.

\subsection{Genuine principal series}

We define the representations and compute their Jacquet modules
directly, so that nothing depends on the general theory of
\cite{Mc12}, whose hypotheses are stronger than ours.

Let $\chi$ be a genuine smooth finite-dimensional representation of
$\At$ on $V_\chi$, extended to $\Bt=\At N$ trivially on the image of
$N$. The \emph{normalized principal series}
$\PI(\chi)=\Ind_{\Bt}^{\Gt}(\chi\otimes\delta_B^{1/2})$ acts by right
translation on locally constant $\phi:\Gt\to V_\chi$ with
\begin{equation}\label{eq:induction}
\phi(bg)=\delta_B(b)^{1/2}\chi(b)\phi(g),\qquad b\in\Bt,\ g\in\Gt .
\end{equation}
It is genuine and admissible. Put
\begin{equation}\label{eq:weyl-action}
\chi^{\wt}(a)=\chi(\wt^{-1}a\wt).
\end{equation}
By Lemma~\ref{lem:weyl-torus}, $\chi^{\wt}$ is not
$a\mapsto\chi(a^{-1})$ in general. For a smooth representation $V$ of
$\Bt$ let $V_N$ be its space of coinvariants under the image of $N$, and
let $r_N(V)$ be $V_N$ with the action of $\At$ twisted by
$\delta_B^{-1/2}$. This is the normalized Jacquet module. We use two
general facts about $\ell$-groups \cite{BZ76,BZ77}. First, $V\mapsto
V_N$ is exact, because $N$ is a union of compact open subgroups. Second,
$C_c^\infty(N)_N\cong\Cc$ through $\phi\mapsto\int_N\phi$. We also use
Frobenius reciprocity
$\Hom_{\Gt}(\pi,\PI(\chi))\cong\Hom_{\At}(r_N(\pi),\chi)$, which holds
for parabolic induction in any $\ell$-group of this form
\cite{BZ77}.

\begin{lemma}[Jacquet module]\label{lem:jacquet}
Let $\chi$ be a genuine smooth finite-dimensional representation of
$\At$. Then $r_N(\PI(\chi))$ has a subrepresentation isomorphic to
$\chi^{\wt}$ with quotient isomorphic to $\chi$.
\end{lemma}

\begin{proof}
Let $V=\PI(\chi)$ and $V_0=\{\phi\in V:\phi(1)=0\}$. Since
$G=B\sqcup BwN$, we have $\Gt=\Bt\sqcup\Bt\wt N$, and $\Bt\backslash\Gt$
is compact. A function in $V_0$ vanishes on $\Bt$ and is locally
constant, so it vanishes near $\Bt$. Hence
$\phi\mapsto(u\mapsto\phi(\wt(u,1)))$ is an isomorphism from $V_0$ onto
$C_c^\infty(N,V_\chi)$, and $\phi\mapsto\phi(1)$ identifies $V/V_0$
with $V_\chi$. In both cases the image of $N$ acts by right
translation.

On $V/V_0$, the element $a\in\At$ acts by $\delta_B(a)^{1/2}\chi(a)$,
so $r_N(V/V_0)\cong\chi$. On $V_0$, write $a=\zeta\hl(t)$. Then
$\wt(u,1)a=(\wt a\wt^{-1})\,\wt\,(a^{-1}(u,1)a)$, and conjugation by $a$
preserves the canonical lift of $N$ (Lemma~\ref{lem:cocycle}(c)). So
$a$ sends $F(u)=\phi(\wt(u,1))$ to
$\delta_B(\wt a\wt^{-1})^{1/2}\chi(\wt a\wt^{-1})F(a^{-1}ua)$. Integrating
over $N$, and using $\int_NF(a^{-1}ua)\,du=\delta_B(a)\int_NF(u)\,du$
and $\delta_B(\wt a\wt^{-1})=\delta_B(a)^{-1}$, we find that $a$ acts on
$(V_0)_N\cong V_\chi$ by $\delta_B(a)^{1/2}\chi(\wt a\wt^{-1})$. After
the twist by $\delta_B^{-1/2}$ this is $\chi(\wt a\wt^{-1})$. Since
$\wt^2$ is central (Lemma~\ref{lem:weyl-torus}(d)),
$\chi(\wt a\wt^{-1})=\chi(\wt^{-1}(\wt^2a\wt^{-2})\wt)=\chi^{\wt}(a)$.
Exactness of $r_N$ gives the statement.
\end{proof} For
$\chi=\chi_\lambda$ we have $\chi^{\wt}\cong\chi_{\lambda^w}$ with
$\lambda^w(z)=\lambda(\wt^{-1}z\wt)$, by
Theorem~\ref{thm:torus}. Call $\chi_\lambda$ \emph{Weyl regular} if
$\lambda^w\ne\lambda$.

\begin{proposition}\label{prop:tempered}
Let $\chi=\chi_\lambda$ with $\lambda$ unitary.
\begin{enumerate}[label=(\alph*)]
\item $\PI(\chi)$ is unitary.
\item If $\chi$ is Weyl regular, then $\PI(\chi)$ is irreducible.
\end{enumerate}
\end{proposition}

\begin{proof}
(a) Let $\langle\ ,\ \rangle$ be a $\At$-invariant inner product on
$V_\chi$. For $\phi_1,\phi_2$ in the induced space, the function
$g\mapsto\langle\phi_1(g),\phi_2(g)\rangle$ transforms on the left
under $\Bt$ by $\delta_B$. So its integral over $\Bt\backslash\Gt$
is a $\Gt$-invariant inner product.

(b) By Frobenius reciprocity and Lemma~\ref{lem:jacquet},
\[
\dim\operatorname{End}_{\Gt}(\PI(\chi))
=\dim\Hom_{\At}(r_N\PI(\chi),\chi)
\le\dim\Hom(\chi,\chi)+\dim\Hom(\chi^{\wt},\chi)=1+0 .
\]
Here $\chi=\chi_\lambda$ and $\chi^{\wt}\cong\chi_{\lambda^w}$ are
irreducible and not isomorphic. Now suppose $W\subset\PI(\chi)$ is a
nonzero proper invariant subspace. The orthogonal complement $W^\perp$
of $W$ in $\PI(\chi)$, for the inner product of (a), is invariant. Each
space of $s(K_m)$-fixed vectors is finite-dimensional, by
admissibility, and the orthogonal projection onto it is $\Pi(e_m)$ for
the normalized genuine idempotent $e_m$ of $s(K_m)$. That projection is
self-adjoint and preserves $W$. Splitting each fixed space orthogonally
therefore gives $\PI(\chi)=W\oplus W^\perp$. The projection onto $W$
is then a $\Gt$-endomorphism that is not a scalar, a contradiction.
\end{proof}

\begin{uvremark}\label{uv:geometric}
\emph{Status: cited only.} The proofs of Lemma~\ref{lem:jacquet} and
Proposition~\ref{prop:tempered} use three general facts about
$\ell$-groups: exactness of coinvariants, $C_c^\infty(N)_N\cong\Cc$, and
Frobenius reciprocity \cite{BZ76,BZ77}. They do not depend on the
cover, but their precise theorem numbers were not checked. Luo
\cite{Lu20} proves the Knapp--Stein dimension theorem for finite
central covers, which sharpens Proposition~\ref{prop:tempered}(b). The
unitary principal series are also tempered. That is standard for linear
groups, and \cite{Li12} develops the tempered theory for covers, but it
was not checked here.
\end{uvremark}

\begin{lemma}[Spherical vectors]\label{lem:Kfixed}
Assume $\gcd(p,2n)=1$. Call $\lambda$ \emph{unramified} if
$\lambda((h(u),1))=1$ for all $u\in(\Of^\times)^d$. Then
$\dim\PI(\chi_\lambda)^{s(K)}$ is $1$ if $\lambda$ is unramified, and
$0$ otherwise.
\end{lemma}

\begin{proof}
Since $\Gt=\Bt s(K)$, the map $\phi\mapsto\phi(1)$ identifies
$\PI(\chi)^{s(K)}$ with the vectors in $V_\chi$ fixed by
$\Bt\cap s(K)=s(B\cap K)$. The image of $N\cap K$ acts trivially, so
these are the vectors fixed by $s(A\cap K)=\{\hl(u):u\in\Of^\times\}$.
This group lies in $\Ct$, since $A\cap K\subset C$.

The subgroup $\Ct$ is normal in $\At$, since it contains the center
and $\At/\ZA$ is abelian. By Mackey's formula, the restriction of
$\chi_\lambda$ to $\Ct$ is the sum of the $d$ conjugates
$\tilde\lambda^{b}$ with $b=\hl(\varpi^j)$, $0\le j<d$. By
Lemma~\ref{lem:commutator},
$\hl(\varpi^j)^{-1}\hl(u)\hl(\varpi^j)=(u,\varpi^j)_n^{-2}\hl(u)
=(\varpi,u)_n^{2j}\hl(u)$, so
$\tilde\lambda^{\hl(\varpi^j)}(\hl(u))=(\varpi,u)_n^{2j}\,
\tilde\lambda(\hl(u))$. On
$\Of^\times$ the character $u\mapsto(\varpi,u)_n^2$ has order exactly
$d$, by (H5). So the restrictions of the $d$ conjugates to
$s(A\cap K)$ run once through the $d$ characters of $\Of^\times$ that
agree with $\lambda$ on $(\Of^\times)^d$. If $\lambda$ is unramified,
exactly one of them is trivial. Otherwise none is.
\end{proof}

\begin{example}\label{ex:Kfixed-three}
Take $n=3$ and $p>3$. Let $\lambda$ be unramified and put
$\alpha=\lambda(\hl(\varpi^3))$. Then $\PI(\chi_\lambda)$ has a unique
spherical line. By (H5), $(\varpi^3,-1)_3=(\varpi,-1)_3^3=1$, so
Lemma~\ref{lem:weyl-torus}(c) gives
$\lambda^w(\hl(\varpi^3))=\lambda(\hl(\varpi^{-3}))=\alpha^{-1}$. If
$\lambda$ is unitary and $\alpha\ne\pm1$, then $\chi_\lambda$ is Weyl
regular and $\PI(\chi_\lambda)$ is irreducible by
Proposition~\ref{prop:tempered}(b).
\end{example}

\subsection{Reducibility}

Reducible unitary principal series occur only for odd $n$. Nothing in
this subsection is proved here.

\begin{uvremark}\label{uv:reducible}
\emph{Status: source checked} (statements and introductions of
\cite{GS16} and \cite{Sz25}). Moen \cite{Mo88} proves irreducibility of
certain unitary principal series of covers of $\SL(2)$ for odd $p$.
Goldberg and Szpruch \cite[Thm.~5.2]{GS16}, assuming $\mu_n\subset F$
and $p\nmid n$, show that a genuine principal series induced from
unitary data is reducible if and only if $n$ is odd and the projection
of the central character of the inducing representation to the $d$th
powers is a nontrivial quadratic character. Szpruch \cite[\S3.1]{Sz25}
explains that these representations are parametrized by the nontrivial
quadratic characters of $F^\times$, and that each is a sum of two
inequivalent irreducible representations. These constituents must be
included in any complete account of the unitary dual. This agrees with
claim~(c) of Table~\ref{tab:status}.
\end{uvremark}

\subsection{Complementary series and endpoints}\label{subsec:complementary}

A \emph{complementary series} representation is an
irreducible unitarizable $\PI(\chi)$ with $\chi$ not unitary. These
representations are not tempered. Fix a unitary genuine character
$\lambda_0$ of $\ZA$ and consider the family
$\lambda_s(\zeta\hl(u))=\lambda_0(\zeta\hl(u))|u|^{s/d}$ for
$u\in(F^\times)^d$, so that $\lambda_s(\hl(t^d))$ behaves like
$|t|^{s}$. An invariant Hermitian form on $\PI(\chi_{\lambda_s})$, when
it exists, is built from a suitably normalized standard intertwining
operator $M(s)$. Its positivity is a substantive statement. It does
not follow from the uniqueness of invariant Hermitian forms, which
only controls a scalar once irreducibility is known.

\paragraph{Converting the parameter.} The 1999 manuscript induces from a
character $\mu$ of $\Ct$ of the form $\mu_0(x)|x|^{s_{\mathrm{old}}}$,
where $x$ is the torus coordinate of $h(x)\in C$. On the common
subgroup $\ZA$, where $x=t^d$, this gives $|t^d|^{s_{\mathrm{old}}}$,
while $\lambda_s$ gives $|t|^{s}$. So, formally,
\begin{equation}\label{eq:s-conversion}
s=d\,s_{\mathrm{old}} .
\end{equation}
This is a comparison of exponents on $\ZA$ only. The two papers use
different subgroups and sections, so \eqref{eq:s-conversion} must be
checked against each source before endpoint values are copied. Under
\eqref{eq:s-conversion} the 1999 endpoint $s_{\mathrm{old}}=1/n$ becomes
$s=d/n$, which is $1$ for odd $n$ and $1/2$ for even $n$.

\begin{uvremark}\label{uv:complementary}
\emph{Status: cited only}, except for Arit\"urk's reducibility points,
which are source checked against his introduction. For the double cover
with $p$ odd, Gelbart and Sally \cite{GS75,Ge76} treat the complementary
series and the endpoint where the even Weil representation appears.
For $n=2$ we have $d=1$, so the two parameters agree and the endpoint is
$s=1/2$. Arit\"urk \cite{Ar80} treats the threefold cover. In his
parameter, the genuine principal series off $\operatorname{Re}s=0$ are
irreducible except at $s=\pm1/3$, and he describes the composition
series there. The 1999 manuscript asserted, for every $n$, that
$0<s_{\mathrm{old}}<1/n$ gives complementary series and that the
endpoint $s_{\mathrm{old}}=1/n$ has a square-integrable
subrepresentation and a non-tempered spherical quotient. It proved
unitarizability by extending the representation to a cover of
$\GL(2,F)$ and citing \cite[\S\S26--27]{FK86}, but the extension was
not proved. It deduced unitarity of the quotient from
\cite[Thm.~II.2.1]{KP84}. These are claims (d) and (e) of
Table~\ref{tab:status}. They are not proved here, they may need the
hypothesis $p\nmid n$, and their endpoints must be converted by
\eqref{eq:s-conversion}.
\end{uvremark}

\subsection{The intertwining-operator approach}\label{subsec:intertwining}

A complementary-series theorem for general $n$ would come from the
positivity of a normalized intertwining operator. We describe the steps
and where the difficulty lies. The standard intertwining operator is
\[
(T_s\varphi)(g)=\int_F\varphi\bigl(\wt^{-1}(u(x),1)\,g\bigr)\,dx .
\]
Where the integral converges, it maps $\PI(\chi_{\lambda_s})$ to
$\PI(\chi_{\lambda_s}^{\wt})$. Following Arit\"urk \cite{Ar80}, one can
write its action in a Kirillov-type model as a $d\times d$ matrix $J_s$.
Three separate steps are then needed.
\begin{enumerate}[label=(\roman*),leftmargin=*]
\item Meromorphic continuation and normalization of $T_s$, with the
  correct Weyl-conjugate inducing datum.
\item Construction of a nondegenerate invariant Hermitian form from
  $T_s$, for a family that is Hermitian self-dual.
\item Positivity of that form, equivalently of the normalized matrix
  $J_s$, on the whole proposed interval.
\end{enumerate}
None of these follows from the corresponding fact for a cover of
$\GL(2,F)$ until one proves that the representations of $\Gt$ extend.
When $d>1$ the matrix $J_s$ is not a scalar, which makes step (iii)
harder. The local coefficient matrices of \cite{GS16,Sz19} are a
natural tool for steps (i) and (ii). For example, when $n=2$ we have
$d=1$, $J_s$ is a scalar function of $s$, and the classical analysis
of \cite{GS75} applies.

\subsection{Supercuspidal representations and completeness}

Completing the unitary dual needs, besides the principal series, the
supercuspidal representations and the unitary subquotients of the
reducible induced representations.

\begin{uvremark}\label{uv:subquotient}
\emph{Status: cited only.} An irreducible admissible genuine representation is
\emph{supercuspidal} if its Jacquet module along $N$ vanishes. By the
analogue of Jacquet's theorem for covers \cite{BZ77,BD84}, every
irreducible admissible genuine representation is either supercuspidal
or a subquotient of some $\PI(\chi)$. The supercuspidal representations
of the double cover are described in \cite{Ma84a,Ma84b}, and those of
the $n$-fold cover with $\gcd(p,2n)=1$ in \cite{Jo99}. These citations
do not by themselves give the unitary dual for all $n$, because one
must still decide which subquotients of the $\PI(\chi)$ are unitary.
For example, for $n=3$ this requires the composition series at
$s=\pm1/3$ from \cite{Ar80}.
\end{uvremark}

%======================================================================
\section{Principal-series characters}\label{sec:characters}

The trace of $\PI(\chi_\lambda)$ can be computed directly from the
induced model, and its character then follows from the Weyl
integration formula. The character involves Weyl conjugation in the
cover, which is not the same as inversion.

\begin{lemma}[Trace of an induced representation]
\label{lem:induced-trace}
Assume $\gcd(p,2n)=1$. Let $\chi$ be a genuine smooth
finite-dimensional representation of $\At$ and $\Pi=\PI(\chi)$. For
$f\in C_c^\infty(\Gt)_\gen$, the operator $\Pi(f)$ has finite rank and
\[
\tr\Pi(f)=\int_{\At}\tr\chi(a)\,\Psi_f(a)\,da ,
\]
with $\Psi_f$ as in \eqref{eq:Psi}.
\end{lemma}

\begin{proof}
\emph{The compact model.} Since $\Gt=\Bt s(K)$, restriction to $s(K)$
identifies the induced space with the space $H$ of locally constant
$\phi:K\to V_\chi$ with $\phi(bk)=\chi(s(b))\phi(k)$ for $b\in B\cap K$.
We regard $H$ as a subspace of $L^2(K,V_\chi)$, using a fixed inner
product on $V_\chi$.

\emph{The kernel.} For $\phi\in H$ and $k\in K$, use the invariance of
$dg$ and \eqref{eq:iwasawa}:
\[
\Pi(f)\phi(k)=\int_{\Gt}f(s(k)^{-1}g)\phi(g)\,dg
=\int_K\mathcal K(k,k')\phi(k')\,dk',
\]
where $\mathcal K(k,k')=\int_{\Bt}f(s(k)^{-1}bs(k'))\delta_B(b)^{1/2}
\chi(b)\,db$ and $db=da\,du$. Choose $m$ such that $f$ is left and
right invariant under $s(K_m)$, up to the genuine factor. Then
$\mathcal K(k,k')$ depends only on the cosets $kK_m$ and $k'K_m$. So
$\mathcal K$ is a finite sum $\sum_{i,j}1_{k_iK_m}(k)\,
1_{k_jK_m}(k')\,M_{ij}$ with $M_{ij}\in\operatorname{End}(V_\chi)$.

\emph{The operator on $L^2$.} Define $T$ on all of $L^2(K,V_\chi)$ by
$T\phi(k)=\int_K\mathcal K(k,k')\phi(k')\,dk'$. By the previous
paragraph, $T$ has finite rank, its image consists of locally constant
functions, and $\tr T=\int_K\tr\mathcal K(k,k)\,dk$. A substitution
$b\mapsto s(b_0)b$ shows that
$\mathcal K(b_0k,k')=\chi(s(b_0))\mathcal K(k,k')$ for $b_0\in B\cap K$.
So the image of $T$ lies in $H$. On $H$, $T$ agrees with $\Pi(f)$.

\emph{The trace.} Let $E$ be the image of $T$, a finite-dimensional
subspace of $H$. Both $T$, on $L^2(K,V_\chi)$, and $\Pi(f)$, on $H$, map
into $E$, and they agree on $E$. The trace of a finite-rank operator
equals the trace of its restriction to any finite-dimensional invariant
subspace that contains its image. So $\tr\Pi(f)=\tr(T|_E)=\tr T$.
Hence
\[
\tr\Pi(f)=\int_K\tr\mathcal K(k,k)\,dk
=\int_{\At}\int_N\tr\chi(a)\,\delta_B(a)^{1/2}f^K(a(u,1))\,du\,da ,
\]
because $\chi$ is trivial on the image of $N$. This is the stated
formula.
\end{proof}

\begin{corollary}[Trace and orbital integrals]\label{cor:trace-integral}
Assume $\gcd(p,2n)=1$. Let $\lambda$ be a genuine smooth character of
$\ZA$ and $\Pi=\PI(\chi_\lambda)$. For $f\in C_c^\infty(\Gt)_\gen$,
\begin{align}
\tr\Pi(f)&=\frac12\int_{\ZA}\bigl(\tr\chi_\lambda(z)
+\tr\chi_\lambda^{\wt}(z)\bigr)\Phi_f(z)\,dz\label{eq:trace-first}\\
&=d\int_{\ZA}\lambda(z)\Phi_f(z)\,dz .\label{eq:trace-second}
\end{align}
\end{corollary}

\begin{proof}
By Theorem~\ref{thm:torus}, $\tr\chi_\lambda$ vanishes off $\ZA$ and
equals $d\lambda$ on $\ZA$. Lemma~\ref{lem:induced-trace} and
Lemma~\ref{lem:abel} give \eqref{eq:trace-second}, since the singular
set has measure zero. The integrand is invariant under $\mu_n$,
because $\lambda$ and $\Phi_f$ transform by inverse characters. The
map $z\mapsto\wt^{-1}z\wt$ preserves the measure on $\ZA$, since it is
$\zeta\hl(u)\mapsto\zeta\hl(u^{-1})$. Together with the Weyl
invariance of $\Phi_f$, this shows that the two terms of
\eqref{eq:trace-first} are equal. So \eqref{eq:trace-first} equals
\eqref{eq:trace-second}.
\end{proof}

\begin{theorem}[Principal-series character]\label{thm:character}
Assume $\gcd(p,2n)=1$. Let $\lambda$ be a genuine smooth character of
$\ZA$, $\chi=\chi_\lambda$ and $\Pi=\PI(\chi)$, which need not be
irreducible. Define $\Theta_\Pi$ on the regular semisimple set of
$\Gt$ by
\begin{equation}\label{eq:char-formula}
\Theta_\Pi(x^{-1}ax)=
\frac{\tr\chi(a)+\tr\chi^{\wt}(a)}{|D(a)|^{1/2}}\qquad
(a\in\At_\reg,\ x\in\Gt),
\end{equation}
and $\Theta_\Pi=0$ on elliptic regular elements. Then $\Theta_\Pi$ is
a well-defined, locally integrable class function, and
$\tr\Pi(f)=\int_{\Gt}\Theta_\Pi(g)f(g)\,dg$ for all
$f\in C_c^\infty(\Gt)_\gen$. In particular $\Theta_\Pi=0$ on
$\At_\reg\setminus\ZA$, and for $z\in(\ZA)_\reg$
\begin{equation}\label{eq:char-Z}
\Theta_\Pi(z)=\frac{d\bigl(\lambda(z)+\lambda(\wt^{-1}z\wt)\bigr)}
{|D(z)|^{1/2}} .
\end{equation}
It satisfies $\Theta_\Pi(\zeta g)=\zeta\Theta_\Pi(g)$.
\end{theorem}

\begin{proof}
\emph{Well defined.} The right side of \eqref{eq:char-formula} is
invariant under conjugation of $a$ by $\At$, since traces are. It is
also invariant under $a\mapsto\wt^{-1}a\wt$, since
$\tr\chi^{\wt}(\wt^{-1}a\wt)=\tr\chi(\wt^{-2}a\wt^2)=\tr\chi(a)$ by
Lemma~\ref{lem:weyl-torus}(d). If $x^{-1}ax=y^{-1}a'y$ with
$a,a'\in\At_\reg$, then $yx^{-1}$ conjugates $a$ to $a'$ and lies in
$\At\cup\wt\At$. So the value does not depend on the choices.
Formula~\eqref{eq:char-Z} and the vanishing off $\ZA$ follow from
\eqref{eq:torus-trace}, because $\wt^{-1}a\wt\in\ZA$ exactly when
$a\in\ZA$.

\emph{The trace.} Let $f\in C_c^\infty(\Gt)_\gen$. The singular
semisimple set has Haar measure zero, so it suffices to integrate over
regular elements, where $\Theta_\Pi$ is defined. The product
$\Theta_\Pi f$ is invariant under $\mu_n$, so it is the pullback of a
function $\bar\phi$ on $G$, and $\int_{\Gt}\Theta_\Pi f=\int_G\bar\phi$
because $\mu_n$ has mass one. Apply \eqref{eq:weyl}. The elliptic term
vanishes because $\Theta_\Pi$ does. For $a=h(t)$ with
$t\notin(F^\times)^d$ the integrand vanishes. For $t\in(F^\times)^d$
let $z\in\ZA$ lie over $h(t)$. Then
$\bar\phi(x^{-1}h(t)x)=\Theta_\Pi(z)f(\tilde x^{-1}z\tilde x)$ for
any lift $\tilde x$ of $x$, and $A\backslash G$ has the same measure
as $\At\backslash\Gt$. So
\[
\int_{\Gt}\Theta_\Pi f=\frac12\int_{\ZA}|D(z)|\Theta_\Pi(z)
\int_{\At\backslash\Gt}f(x^{-1}zx)\,d\dot x\,dz
=\frac12\int_{\ZA}\bigl(\tr\chi(z)+\tr\chi^{\wt}(z)\bigr)
\Phi_f(z)\,dz,
\]
where we used that the integrand is invariant under $\mu_n$. By
Corollary~\ref{cor:trace-integral} this is $\tr\Pi(f)$.

\emph{Local integrability.} Let $\Omega\subset\Gt$ be compact. We may
assume that $\Omega$ is open and stable under $\mu_n$, so that
$\omega=\rho(\Omega)$ is compact open and $1_\Omega$ is the pullback of
$1_\omega\in C_c^\infty(G)$. The function $|\Theta_\Pi|$ is invariant
under $\mu_n$. The same computation as above, with $|\Theta_\Pi|1_\Omega$
in place of $\Theta_\Pi f$, gives
\[
\int_\Omega|\Theta_\Pi|=\frac12\int_{(F^\times)^d}
\bigl|\tr\chi(z_t)+\tr\chi^{\wt}(z_t)\bigr|\,F^A_{1_\omega}(t)\,dt,
\]
where $z_t\in\ZA$ is any lift of $h(t)$. By \eqref{eq:abel-linear},
$F^A_{1_\omega}$ is bounded and compactly supported. The traces are
continuous, so they are bounded on the lifts of that support. Hence the
integral is finite.

\emph{Equivariance.} Both $\chi(\zeta a)$ and $\chi^{\wt}(\zeta a)$ are
$\zeta$ times their values at $a$.
\end{proof}

\begin{remark}\label{rem:reducible-character}
$\Theta_\Pi$ is the character of the induced representation, which may
be reducible. When $\PI(\chi_\lambda)$ is reducible, $\Theta_\Pi$ is the
sum of the characters of its constituents. The paper does not compute
the characters of the individual constituents.
\end{remark}

\begin{warning}[Inversion is not Weyl conjugation]\label{caution:inverse}
For $n>2$ the expression $\tr\chi(a)+\tr\chi(a^{-1})$ transforms under
$a\mapsto\zeta a$ by $\zeta$ in the first term and by $\zeta^{-1}$ in
the second. Its quotient by $|D(a)|^{1/2}$ is therefore not the
character of a genuine representation, and it cannot be used in a
matching theorem. By Lemma~\ref{lem:weyl-torus}(b), Weyl conjugation
looks like inversion on the matrix part. It differs on the central
part, and Example~\ref{ex:inverse} shows that it can differ even on
$\hl(t)$.
\end{warning}

\begin{example}\label{ex:char-three}
Take $n=3$ and $p>3$, and let $\lambda$ be a genuine character of $\ZA$
with $\alpha=\lambda(\hl(\varpi^3))$. As in
Example~\ref{ex:Kfixed-three}, $\lambda(\wt^{-1}\hl(\varpi^3)\wt)=
\lambda(\hl(\varpi^{-3}))=\alpha^{-1}$. Since
$|\varpi^3-\varpi^{-3}|=q^3$, Theorem~\ref{thm:character} gives
\[
\Theta_\Pi(\hl(\varpi^3))=\frac{3(\alpha+\alpha^{-1})}{q^3},\qquad
\Theta_\Pi(\hl(\varpi))=0 .
\]
The second value vanishes because $\varpi\notin(F^\times)^3$.
\end{example}

%======================================================================
\section{A restricted and normalized split-torus transfer}
\label{sec:transfer}

Genuine principal series of $\Gt$ can be compared with principal
series of $G$ through a power map from $A$ to $\ZA$. That map is a
homomorphism only for some $n$. In this section assume
$\gcd(p,2n)=1$ and
\begin{equation}\label{eq:restriction}
n\text{ is odd}\qquad\text{or}\qquad 4\mid n .
\end{equation}
Then $n\mid d^2$. Indeed $d=n$ when $n$ is odd, and $d=n/2$ is even
when $4\mid n$. For $t\in F^\times$ define
\begin{equation}\label{eq:nu}
\nu(h(t))=(h(t^d),1)\in\ZA .
\end{equation}

\begin{lemma}\label{lem:norm}
Under \eqref{eq:restriction}, $\nu:A\to\ZA$ is a continuous
homomorphism with kernel $h(\mu_d)$, of order $d$. Its image is
$\ZA^0$, which is a subgroup, and $\ZA=\mu_n\times\ZA^0$. Moreover
\begin{equation}\label{eq:nu-w}
\wt^{-1}\nu(h(t))\wt=\nu(h(t^{-1})).
\end{equation}
If $dz$ is the restriction to $\ZA^0$ of the measure on $\At$, then
\begin{equation}\label{eq:haar-push}
\int_{\ZA^0}H(z)\,dz=\frac1d\int_AH(\nu(h(t)))\,dt
\end{equation}
for every compactly supported locally constant $H$ on $\ZA^0$.
\end{lemma}

\begin{proof}
By Lemma~\ref{lem:cocycle}(a),
$\beta(h(s^d),h(t^d))=(t^d,s^d)_n=(t,s)_n^{d^2}=1$, since $n\mid d^2$.
So $\nu$ is a homomorphism, its image $\ZA^0$ is a subgroup, and
$\ZA^0\cap\mu_n=1$. The kernel is $\{h(t):t^d=1\}$, which has order
$d$ because $d\mid n$ and $\mu_n\subset F$. Equation~\eqref{eq:nu-w} is
Lemma~\ref{lem:weyl-torus}(b) with $t^d$ in place of $t$.

For \eqref{eq:haar-push}, the $d$th power map sends $\varpi^k\Of^\times$
onto $\varpi^{dk}(\Of^\times)^d$, and every element of $(F^\times)^d$
lies in exactly one of these sets. The first set has volume $1$, and
the second has restricted volume $1/d$ by Example~\ref{ex:volume}. So
the pushforward of $dt$ under $t\mapsto t^d$ is $d$ times the restricted
measure on $(F^\times)^d$. This is \eqref{eq:haar-push}.
\end{proof}

\begin{definition}[Matching on the split torus]\label{def:match}
Let $\bar f\in C_c^\infty(\Gt)_\gen$ and $f\in C_c^\infty(G)$. They are
\emph{split-matching} if
\begin{equation}\label{eq:match}
F_f^A(t)=\Phi_{\bar f}(\nu(h(t)))\qquad(t\in F^\times),
\end{equation}
where both sides are the locally constant extensions of
Sections~\ref{subsec:orbital-linear} and~\ref{sec:orbital}. This
notion concerns split orbital integrals only.
\end{definition}

\begin{example}\label{ex:units-match}
The genuine unit $e_K$ and $1_K$ are split-matching. Indeed
$\Phi_{e_K}(\nu(h(t)))=1_{\Of^\times}(t^d)=1_{\Of^\times}(t)$ by
Example~\ref{ex:Phi-unit}, and $F_{1_K}^A(t)=1_{\Of^\times}(t)$ by
Example~\ref{ex:FA-unit}.
\end{example}

\begin{theorem}[Restricted split-torus trace matching]
\label{thm:transfer}
Assume $\gcd(p,2n)=1$ and \eqref{eq:restriction}. For a genuine smooth
character $\lambda$ of $\ZA$ put
\begin{equation}\label{eq:eta}
\eta_\lambda(t)=\lambda(\nu(h(t))),\qquad t\in F^\times .
\end{equation}
This is a smooth character of $A\cong F^\times$, unitary if $\lambda$
is. If $\bar f$ and $f$ are split-matching, then
\begin{equation}\label{eq:transfer-identity}
\boxed{\ \tr\PI(\chi_\lambda)(\bar f)=
\tr\Ind_B^G(\eta_\lambda\delta_B^{1/2})(f)\ }
\end{equation}
Every $\bar f\in C_c^\infty(\Gt)_\gen$ has at least one split-matching
$f\in C_c^\infty(G)$. Neither induced representation needs to be
irreducible.
\end{theorem}

Theorem~\ref{thm:transfer} is an equality of traces for split-matching
test functions, and nothing more. It does not give a transfer operator,
a matching of elliptic orbital integrals, a correspondence between all
irreducible representations, or identities for the two constituents of
a reducible principal series.

\begin{proof}
By Lemma~\ref{lem:norm}, $\eta_\lambda$ is a character. The integrand
$\lambda\Phi_{\bar f}$ is invariant under $\mu_n$, and $\mu_n$ has mass
one. By Corollary~\ref{cor:trace-integral} and Lemma~\ref{lem:norm},
\[
\tr\PI(\chi_\lambda)(\bar f)
=d\int_{\ZA^0}\lambda(z)\Phi_{\bar f}(z)\,dz
=\int_A\eta_\lambda(t)\Phi_{\bar f}(\nu(h(t)))\,dt
=\int_A\eta_\lambda(t)F_f^A(t)\,dt .
\]
By \eqref{eq:trace-linear} the last integral is
$\tr\Ind_B^G(\eta_\lambda\delta_B^{1/2})(f)$.

For existence, $\Phi_{\bar f}\in C_c^\infty(\ZA)_\gen^W$ by
Theorem~\ref{thm:orbital}. The map $\nu$ is proper, since its kernel is
finite and its image is closed. So $t\mapsto\Phi_{\bar f}(\nu(h(t)))$
is locally constant and compactly supported, and it is invariant under
$t\mapsto t^{-1}$ by \eqref{eq:nu-w}. The case $n=1$ of
Theorem~\ref{thm:orbital} gives $f\in C_c^\infty(G)$ with this
function as $F_f^A$.
\end{proof}

\begin{example}[The normalization in degree three]\label{ex:degree-three}
Take $n=3$ and $p>3$, and let $\lambda$ be unramified. By
Example~\ref{ex:Phi-unit} and Example~\ref{ex:volume},
\[
\tr\PI(\chi_\lambda)(e_K)=3\int_{\ZA}\lambda(z)\Phi_{e_K}(z)\,dz
=3\vol\bigl((\Of^\times)^3\bigr)=3\cdot\tfrac13=1 .
\]
This agrees with Lemma~\ref{lem:Kfixed} and Example~\ref{ex:eK}. It
also agrees with Theorem~\ref{thm:transfer}, since $e_K$ matches $1_K$
(Example~\ref{ex:units-match}) and the right side is $1$ by
Example~\ref{ex:trace-linear}. The example shows why the factor in
\eqref{eq:trace-second} is $d$. Using the index $d^2$ there would give
$3$ instead of $1$.
\end{example}

\begin{remark}[The case $n\equiv2\pmod4$]\label{rem:weil}
If $n\equiv2\pmod4$, then $d$ is odd and $n\nmid d^2$. The map
$t\mapsto(h(t^d),1)$ is then not a homomorphism in general, and
$\ZA^0$ is not a subgroup. The phrase ``twist by the Weil index'' does
not by itself define a $\mu_n$-valued homomorphism here. Weil indices
take values in $\mu_4$ and need a compatible normalization or a larger
central group. So \eqref{eq:transfer-identity} is not claimed for
$n\equiv2\pmod4$, and in particular not for $n=2$. This is a limitation
of the particular power map $\nu$, not a statement that no transfer
exists in these degrees. The double cover must be treated with its own
genuine characters and transfer factors \cite{Ge76,GS75,Fl80,Sc98}.
\end{remark}

\begin{remark}[Spherical functions]\label{rem:spherical}
The identity \eqref{eq:transfer-identity} holds whenever $\bar f$ and
$f$ are spherical and split-matching. It is not proved here that every
spherical $\bar f$ has a spherical partner. That is a statement about
the metaplectic Satake transform. Example~\ref{ex:units-match} shows
that it holds for the units.
\end{remark}

%======================================================================
\section{Summary of results and open questions}\label{sec:limitations}

We collect what is proved, what is cited, a running example, and the
open questions.

\paragraph{Proved here.} Lemma~\ref{lem:radical} and
Theorem~\ref{thm:torus} describe the covering torus. They distinguish
three integers that are easily confused:
\[
\dim\chi_\lambda=d,\qquad[\At:\ZA]=d^2,\qquad|\ker\nu|=d .
\]
Lemma~\ref{lem:vanishing} and Theorem~\ref{thm:orbital} give the
vanishing and the image of split orbital integrals.
Lemma~\ref{lem:induced-trace}, Corollary~\ref{cor:trace-integral} and
Theorem~\ref{thm:character} give the trace and the character of
genuine principal series. Proposition~\ref{prop:tempered} and
Lemma~\ref{lem:Kfixed} give unitarity, irreducibility for Weyl regular
data and the spherical vectors. Theorem~\ref{thm:transfer} gives the
split-torus character identity for $\gcd(p,2n)=1$ and $n$ odd or
divisible by $4$. For instance, for $n=3$ and $F=\Qq_7$ it compares
$\PI(\chi_\lambda)$ with $\Ind_B^G(\eta_\lambda\delta_B^{1/2})$, where
$\eta_\lambda(t)=\lambda(\hl(t^3))$.

\paragraph{Not proved here.} Remarks~\ref{uv:geometric} (cited only),
\ref{uv:reducible} (source checked), \ref{uv:complementary} (mostly
cited only) and~\ref{uv:subquotient} (cited only). Claims (d), (e) and
(f) of Table~\ref{tab:status} are neither proved nor refuted.

\paragraph{A running example.} Table~\ref{tab:running} collects the
threefold cover over $F=\Qq_7$, with $\varpi=7$ and $q=7$. Here $d=3$.

\begin{table}[ht]
\centering
\small
\renewcommand{\arraystretch}{1.15}
\begin{tabular}{@{}lll@{}}
\toprule
Quantity & Value & Where\\
\midrule
$(7,3)_3$ & image of $4\bmod 7$, a primitive cube root of $1$ &
Example~\ref{ex:hilbert}\\
$\hl(7)\hl(3)$ versus $\hl(3)\hl(7)$ & differ by a nontrivial cube root
of $1$ & Example~\ref{ex:noncommute}\\
$[\At:\ZA]$, $[\At:\Ct]$, $\dim\chi_\lambda$ & $9$, $3$, $3$ &
Lemma~\ref{lem:radical}, Example~\ref{ex:two-three}\\
$\mathcal O_f^{\ZA}(\hl(7))$ & $0$, since $7\notin(\Qq_7^\times)^3$ &
Example~\ref{ex:vanishing}\\
$\vol\bigl((\Zz_7^\times)^3\bigr)$ & $1/3$ & Example~\ref{ex:volume}\\
$\Phi_{e_K}(\zeta\hl(u))$, $u\in(\Qq_7^\times)^3$ &
$\zeta^{-1}1_{\Zz_7^\times}(u)$ & Example~\ref{ex:Phi-unit}\\
$\dim\PI(\chi_\lambda)^{s(K)}$, $\lambda$ unramified & $1$ &
Lemma~\ref{lem:Kfixed}\\
$\Theta_\Pi(\hl(7^3))$, $\Theta_\Pi(\hl(7))$ &
$3(\alpha+\alpha^{-1})/7^3$, $0$ & Example~\ref{ex:char-three}\\
$\tr\PI(\chi_\lambda)(e_K)$ & $1=\tr\Ind_B^G(\eta_\lambda\delta_B^{1/2})(1_K)$
& Example~\ref{ex:degree-three}\\
$\eta_\lambda(t)$ & $\lambda(\hl(t^3))$ & Theorem~\ref{thm:transfer}\\
\bottomrule
\end{tabular}
\caption{The threefold cover of $\SL(2,\Qq_7)$. Here
$\alpha=\lambda(\hl(7^3))$.}
\label{tab:running}
\end{table}

\paragraph{Open questions.} These are logically independent.
\begin{enumerate}[label=(\arabic*),leftmargin=*]
\item \emph{The 1999 and 2001 unitary-dual claims.} For
  $\gcd(p,2n)=1$, independently prove or refute claims (d)--(f) of
  Table~\ref{tab:status}: the complementary intervals, the constituents
  at their endpoints, and the unitarity of those constituents.
\item \emph{Wild residue characteristic.} Treat $p\mid n$. The center is
  still $\ZA$ by Lemma~\ref{lem:radical}, but the index $[\At:\ZA]$ and
  the splitting of $K$ change. The work of Szpruch \cite{Sz19,Sz25} does
  not assume $p\nmid n$.
\item \emph{Degrees $n\equiv2\pmod4$.} The power map $\nu$ fails here.
  The map $t\mapsto(h(t^{2d}),1)$ is a homomorphism, since
  $(t,s)_n^{4d^2}=1$, but its image is smaller than $\ZA^0$. Find a
  transfer in these degrees, possibly with a Weil-index normalization.
\item \emph{Spherical matching.} Decide whether every spherical genuine
  function has a spherical split-matching partner, and relate this to
  the Satake isomorphism of \cite{Mc12}.
\item \emph{Elliptic matching and constituents.} Extend the comparison
  to elliptic orbital integrals, and find character identities for the
  constituents of reducible principal series. For $n=2$ this was the
  goal of the 1999 manuscript, as a step toward Labesse's multiplicity
  one conjecture.
\end{enumerate}

%======================================================================
\section{Glossary}\label{sec:glossary}

\begin{description}[leftmargin=3.5cm,style=nextline]
\item[Alternating pairing] A pairing $\langle a,b\rangle$ with
  $\langle a,a\rangle=1$. The commutator on $\At/\ZA$ is one.
\item[Center of $\At$] The group $\ZA=\rho^{-1}(h((F^\times)^d))$. It
  is the radical of the commutator pairing.
\item[Complementary series] Irreducible unitarizable principal series
  induced from a nonunitary character. They are not tempered.
\item[Genuine] A representation on which $\zeta\in\mu_n$ acts by
  $\zeta$. A test function with $f(\zeta g)=\zeta^{-1}f(g)$.
\item[Genuine unit] The function $e_K$ with $e_K(\zeta s(k))=\zeta^{-1}$.
  It acts as the projection onto spherical vectors.
\item[Jacquet module] The coinvariants of a representation under $N$,
  twisted by $\delta_B^{-1/2}$. For $\PI(\chi)$ it has constituents
  $\chi$ and $\chi^{\wt}$.
\item[Hilbert symbol] The pairing $(\ ,\ )_n:F^\times\times F^\times
  \to\mu_n$ of local class field theory.
\item[Kubota cocycle] The two-cocycle \eqref{eq:kubota} that defines
  $\Gt$.
\item[Normalized orbital integral] An orbital integral multiplied by
  $|D|^{1/2}$. Its quotient group must centralize the element.
\item[Normalized induction] Induction from $\chi\otimes\delta_B^{1/2}$.
  It preserves unitarity.
\item[Power map] The map $\nu(h(t))=(h(t^d),1)$. It is a homomorphism
  when $n$ is odd or $4\mid n$.
\item[Principal series] A representation $\PI(\chi)$ induced from a
  genuine irreducible representation of $\At$.
\item[Spherical] Fixed by $s(K)$, for a fixed splitting $s$.
\item[Split-matching] Equality \eqref{eq:match} of split orbital
  integrals. It is not a stable or endoscopic transfer.
\item[Stone--von Neumann theorem] The parametrization of irreducible
  representations of a two-step nilpotent group by central characters.
\item[Supercuspidal] An irreducible admissible representation whose
  Jacquet module along $N$ vanishes.
\item[Tame] The case $p\nmid n$, where the Hilbert symbol is given by
  (H5).
\item[Weyl conjugate] The representation
  $\chi^{\wt}(a)=\chi(\wt^{-1}a\wt)$. It is not $\chi(a^{-1})$.
\item[Weyl discriminant] $D(h(t))=-(t-t^{-1})^2$.
\item[Weyl regular] $\chi_\lambda$ with $\lambda^w\ne\lambda$.
\end{description}

%======================================================================
\section*{Acknowledgments}
\addcontentsline{toc}{section}{Acknowledgments}

This revision was carried out with three AI systems. Argon-Gemini
refereed the 1999 manuscript, suggested corrections, added a glossary
and produced a first revised manuscript. A later attempt to continue
the revision with Argon-Gemini was not completed. Sol-ChatGPT (OpenAI)
wrote an independent referee report and prepared the corrected
edition that is Draft~1 of this revision. Claude (Anthropic) prepared
Draft~2. Claude rewrote Draft~1 to follow the author's style guide,
checked its mathematics independently, and supplied complete proofs
where Draft~1 gave sketches. These are Lemma~\ref{lem:abel},
Theorem~\ref{thm:orbital}, Proposition~\ref{prop:tempered},
Lemma~\ref{lem:Kfixed}, Lemma~\ref{lem:induced-trace} and
Theorem~\ref{thm:character}. Claude also wrote
Section~\ref{sec:background} and the examples, expanded the literature
review, checked the bibliographic data, and wrote a verification
script. Sol-ChatGPT then refereed Draft~2. Claude prepared this
Draft~3 in response. It added the proofs of Lemma~\ref{lem:splitting}
and Lemma~\ref{lem:jacquet}, Table~\ref{tab:status} on the 1999 and
2001 claims, the parameter conversion \eqref{eq:s-conversion}, and
Table~\ref{tab:running}. AI review is not an independent
certification of the results. Further details on the AI usage are in
the revision log that accompanies this draft.

I thank Courtney Moen for helpful criticisms of the 1999 manuscript,
and an anonymous referee of that manuscript for several useful
suggestions.

%======================================================================
\begin{thebibliography}{Ma84b}

\bibitem[Ar80]{Ar80} H.~Arit\"urk, On the composition series of
principal series representations of a three-fold covering group of
$SL(2,K)$, \emph{Nagoya Math. J.} \textbf{77} (1980), 177--196.

\bibitem[BD84]{BD84} J.~Bernstein, Le ``centre'' de Bernstein,
r\'edig\'e par P.~Deligne, in \emph{Repr\'esentations des groupes
r\'eductifs sur un corps local}, Travaux en Cours, Hermann, Paris,
1984.

\bibitem[BZ76]{BZ76} I.~N.~Bernstein and A.~V.~Zelevinsky,
Representations of the group $GL(n,F)$, where $F$ is a
non-archimedean local field, \emph{Russian Math. Surveys}
\textbf{31} (1976), 1--68.

\bibitem[BZ77]{BZ77} I.~N.~Bernstein and A.~V.~Zelevinsky, Induced
representations of reductive $p$-adic groups.~I, \emph{Ann. Sci.
\'Ecole Norm. Sup.} \textbf{10} (1977), 441--472.

\bibitem[FK86]{FK86} Y.~Flicker and D.~Kazhdan, Metaplectic
correspondence, \emph{Publ. Math. IH\'ES} \textbf{64} (1986), 53--110.

\bibitem[Fl80]{Fl80} Y.~Flicker, Automorphic forms on covering groups
of $GL(2)$, \emph{Invent. Math.} \textbf{57} (1980), 119--182.

\bibitem[Ge76]{Ge76} S.~Gelbart, \emph{Weil's Representation and the
Spectrum of the Metaplectic Group}, Lecture Notes in Math.~530,
Springer, 1976.

\bibitem[GS75]{GS75} S.~Gelbart and P.~J.~Sally, Intertwining operators
and automorphic forms for the metaplectic group, \emph{Proc. Nat.
Acad. Sci. USA} \textbf{72} (1975), 1406--1410.

\bibitem[GS16]{GS16} D.~Goldberg and D.~Szpruch, Plancherel measures
for coverings of $p$-adic $SL_2(F)$, \emph{Int. J. Number Theory}
\textbf{12} (2016), no.~7, 1907--1936.

\bibitem[Jo88]{Jo88} D.~Joyner, On the metaplectic analog of Kazhdan's
``endoscopic'' lifting, \emph{Israel J. Math.} \textbf{61} (1988),
113--154.

\bibitem[Jo98]{Jo98} D.~Joyner, A correspondence for the generalized
Hecke algebra of the metaplectic cover $\overline{SL(2,F)}$, $F$
$p$-adic, \emph{New York J. Math.} \textbf{4} (1998), 223--235.

\bibitem[Jo99]{Jo99} D.~Joyner, A correspondence for the supercuspidal
representations of $\overline{SL_2(F)}$, $F$ $p$-adic, \emph{Arch.
Math.} \textbf{73} (1999), 332--340.

\bibitem[Jo01]{Jo01} D.~Joyner, Invariant distributions on the $n$-fold
metaplectic covers of $GL(r,F)$, $F$ $p$-adic, \emph{J. Fourier Anal.
Appl.} \textbf{7} (2001), no.~4, 343--358.

\bibitem[KP84]{KP84} D.~Kazhdan and S.~J.~Patterson, Metaplectic forms,
\emph{Publ. Math. IH\'ES} \textbf{59} (1984), 35--142.

\bibitem[Ku67]{Ku67} T.~Kubota, Topological covering of $SL(2)$ over a
local field, \emph{J. Math. Soc. Japan} \textbf{19} (1967), 114--121.

\bibitem[Li12]{Li12} W.-W.~Li, La formule des traces pour les
rev\^etements de groupes r\'eductifs connexes. II. Analyse harmonique
locale, \emph{Ann. Sci. \'Ec. Norm. Sup\'er.} \textbf{45} (2012),
787--859.

\bibitem[Lu20]{Lu20} C.~Luo, Knapp--Stein dimension theorem for finite
central covering groups, \emph{Pacific J. Math.} \textbf{306} (2020),
no.~1, 265--280.

\bibitem[Ma84a]{Ma84a} D.~Manderscheid, On the supercuspidal
representations of $SL_2$ and its two-fold cover.~I, \emph{Math.
Ann.} \textbf{266} (1984), 287--296.

\bibitem[Ma84b]{Ma84b} D.~Manderscheid, On the supercuspidal
representations of $SL_2$ and its two-fold cover.~II, \emph{Math.
Ann.} \textbf{266} (1984), 297--306.

\bibitem[Mc12]{Mc12} P.~J.~McNamara, Principal series representations
of metaplectic groups over local fields, in \emph{Multiple Dirichlet
Series, L-functions and Automorphic Forms}, Progr. Math.~300,
Birkh\"auser/Springer, New York, 2012, 299--327. Preprint version
arXiv:0911.2274.

\bibitem[Mo]{Mo} C.~Moen, Intertwining operators for covering groups of
$SL(2)$ over a local field, preprint, undated.

\bibitem[Mo88]{Mo88} C.~H.~Moen, Irreducibility of unitary principal
series for covering groups of $SL(2,k)$, \emph{Pacific J. Math.}
\textbf{135} (1988), no.~1, 89--110.

\bibitem[Ne86]{Ne86} J.~Neukirch, \emph{Class Field Theory},
Grundlehren der Math. Wiss.~280, Springer, 1986.

\bibitem[Sc98]{Sc98} J.~P.~Schultz, \emph{Lifting of Characters of
$\widetilde{SL}_2(F)$ and $SO_{1,2}(F)$ for $F$ a Nonarchimedean Local
Field}, Ph.D. thesis, University of Maryland, 1998.

\bibitem[Si79]{Si79} A.~J.~Silberger, \emph{Introduction to Harmonic
Analysis on Reductive $p$-adic Groups}, Mathematical Notes~23,
Princeton University Press, 1979.

\bibitem[Sz19]{Sz19} D.~Szpruch, On Shahidi local coefficients matrix,
\emph{Manuscripta Math.} \textbf{159} (2019), no.~1--2, 117--159.

\bibitem[Sz25]{Sz25} D.~Szpruch, Whittaker spaces for reducible unitary
principal series representations of coverings of $SL(2,F)$,
\emph{Proc. Amer. Math. Soc.} \textbf{153} (2025), no.~5, 1933--1945.
Preprint version arXiv:2411.11460.

\bibitem[Ta17]{Ta17} S.~Tang, Principal series representations of
metaplectic groups, arXiv:1706.05145, 2017.

\end{thebibliography}

\end{document}
